\section{SL(2,C) vertex data used here}
\label{ch:sl2c-vertex}

This appendix collects the computations on the $\group{SL}(2, \C)$ quartic vertex~\cite{Erbin:2023:SL2CQuarticVertex} that the level-zero statements of \cref{ch:momentum-space} use: the tile integral $J(p, q)$ from which the off-shell four-tachyon contact term is built, the reduction of that contact term to the Mandelstam invariants, the mapping radius at the puncture at infinity, the asymptotics of $J$ at the corner of its region, and the invariance of the contact region under complex conjugation, which is the geometric input of the reality of the quartic vertex (\cref{ch:conventions}).
None of this is in the paper; \cref{ch:momentum-space} uses only the resulting form~\eqref{eq:mom:F} of the contact term.
The first part is a note on $J$, reproduced in full in the notation of~\cite{Erbin:2023:SL2CQuarticVertex}; the four sections after it are written for this document from the derivations behind them.
Two properties of the explicit vertex family at its declared section, the closed form of the mapping radii and the assembly of the contact term from twelve tiles, come from the construction of the family in~\cite{Erbin:2026:SL2CVerticesClosed}; they are taken here as inputs and stated where they are used.
The numerical values of the vertex family itself (its overall normalization and the couplings of the tachyon and the dilaton) are not reproduced here, and are to be given in~\cite{Erbin:2026:SL2CVerticesClosed}.

\paragraph{The tile integral.}
The first block is the note on the integral
\begin{equation*}
  J(p, q)
    = \int_R \abs{Z}^p \abs{1 - Z}^q \, \dd^2 Z
\end{equation*}
over the fundamental domain $R$ of the vertex region, for real $(p, q)$ and $\beta \ge \beta_{04} = \frac{1}{2} + \frac{1}{\sqrt{2}}$, where the region has four boundary arcs.
Its results have different statuses.\record{047-j-integral-closed-form (R50)}
\begin{itemize}
  \item The boundary representation $J = \frac{1}{2} \Im \oint_{\partial R} \overline{G} \, \dd G$, with $G$ an incomplete Beta function; the constant-phase criterion, by which an arc has a straight image under $G$, and then closes in endpoint values of $G$, exactly when the phase of $f \, \dd Z$ is constant on it, with the four loci it gives; the closed forms on the lines $p = 0$ and $q = 0$ and at the point $(-1, -1)$; the reduction of a generic curved arc to a Kampé de Fériet double series; the confocal trapezoid; and the $n^{-3}$ convergence of the nested quadrature: \status{Proven}.
  \item That no closed form is reachable by the boundary route for generic $(p, q)$, off the lines $p = 0$, $q = 0$ and the point $(-1, -1)$: \status{Argued}.
    This is a statement about one route, within the class of finite combinations of ${}_2F_1$ and Appell $F_1$, $F_2$ functions, not a proof that no closed form exists; the step on the ellipse arc, where the ${}_2F_1$ sits at a quadratic-rational argument, is an argument and was not machine-checked.
  \item The evaluation by Chebyshev cumulative integration along $\partial R$ and its accuracy: \status{Numerical}, at $\abs{p} + \abs{q} \le 170$ and $\beta \in \{ \frac{3}{2}, 2 \}$, without a rigorous error bound.
\end{itemize}
The relevance of $J$ to the contact term, and to the pair kernel of \cref{ch:momentum-space}, rests on the two inputs of the vertex family named above (\cref{sec:sl2c:stu}).

\begin{quote}
  \small
  \textbf{Translation box: tile-integral note $\to$ this document.}
  \begin{itemize}
    \setlength{\itemsep}{0.3ex}
    \item The note is written in the notation of~\cite{Erbin:2023:SL2CQuarticVertex}: $Z \equiv Z_4$ is the cross-ratio, $\beta$ the parameter of the family, $C_1$--$C_4$ the four boundary arcs, $a$ and $b$ the semi-axes of the $s$-channel ellipse; $R = \mathrm{FD} \cap \mc V_{0,4}$ is the region written $R$ in~\eqref{eq:mom:F}.
    \item The measure is $\dd^2 Z = \dd x \, \dd y$ for $Z = x + \I y$, as in~\eqref{eq:mom:F} and in the sections below; the normalization of the measure only rescales the overall constant of the contact term.
    \item The Mandelstam invariants are those of the mostly-plus metric with all momenta incoming, $s = (k_1 + k_2)^2$, $t = (k_1 + k_4)^2$, $u = (k_1 + k_3)^2$, with $s + t + u = \sum_a k_a^2$ and the tachyon shell at $\alpha' k_a^2 = 4$.
      In the sections below the invariants are written as $s_{ab} = (k_a + k_b)^2$, so that $t = s_{14} = s_{23}$ and $u = s_{13} = s_{24}$.
    \item In the section on numerics, $p$ and $q$ are the exponents of $J$; they are unrelated to the relative momenta $q, q'$ of \cref{ch:momentum-space}.
    \item The note contains no string field. The overall normalization of the contact term, which differs by a factor $\frac{1}{2}$ per sphere correlator between the normalization of~\cite{Erbin:2023:SL2CQuarticVertex} and the one of this document (\cref{sec:conv:two-normalizations}), does not enter $J$.
    \item \emph{Reality.} The conjugation $\ddagger$ of the section on conjugation invariance below is the one of the earlier notes: on a $k$-string tensor whose strands have ghost number two, $\ddagger = (-1)^k \star$, so that the conditions agree for $\Sigma_2$ and differ by a sign for odd $k$, $\Sigma_3^\star = - \Sigma_3$ and $\Psi^\star = - \Psi$ (\cref{ch:conventions}).
  \end{itemize}
\end{quote}

\subsection{Introduction}
\label{jb:sec:intro}


An off-shell four-point amplitude computed from the $\group{SL}(2, \C)$ quartic vertex of
\cite{Erbin:2023:SL2CQuarticVertex} ends with an integral over the fundamental domain of the vertex
region $\mc V_{0,4}$ in the moduli space of the four-punctured sphere.
For four tachyons the integrand is a power of the distance to each of the two punctures at $0$ and
$1$, as in the Virasoro--Shapiro integral, so what has to be computed is
\begin{equation}
  \label{jb:eq:J-def}
  J(p, q)
    := \int_{R} \abs{Z}^{p} \, \abs{1 - Z}^{q} \, \dd^2 Z,
\end{equation}
with $R$ the fundamental domain of the vertex region and $\dd^2 Z = \dd x \, \dd y$ for
$Z = x + \I y$.
Throughout we abbreviate $Z \equiv Z_4$ for the cross-ratio coordinate on $\mc M_{0,4}$.
The exponents $p$ and $q$ are linear in the Mandelstam invariants: in the application to the cubic
reformulation of the closed-string action (\cref{sec:sl2c:stu}), the off-shell four-tachyon contact term is
a sum of six copies of \eqref{jb:eq:J-def}, one per ordering $(x, y, z)$ of $(s, t, u)$, at
\begin{equation}
  \label{jb:eq:pq-mandelstam}
  (p, q)
    = \frac{\alpha'}{4} \, (x - y - z, \ y - x - z).
\end{equation}
The invariants are those of the mostly-plus metric with all momenta incoming, $s = (k_1 + k_2)^2$,
$t = (k_1 + k_4)^2$, $u = (k_1 + k_3)^2$, so that $s + t + u = \sum_a k_a^2$ and the tachyon shell
is $\alpha' k_a^2 = 4$; a reader using the opposite sign for the invariants should translate
\eqref{jb:eq:pq-mandelstam} accordingly.
A single amplitude point therefore costs six evaluations of \eqref{jb:eq:J-def}, and a scan over
momenta costs thousands.
Done by two-dimensional quadrature this is slow and, worse, not very accurate: the exclusion of the
unit disk gives the natural nested quadrature a square-root branch point in its inner limits, and
the convergence is only algebraic (\cref{jb:sec:numerics}).

This note establishes that \eqref{jb:eq:J-def} is exactly a one-dimensional integral along
$\partial R$, and analyzes how far that reduction can be pushed towards a closed form.
The result is a property of the vertex region of \cite{Erbin:2023:SL2CQuarticVertex}, not of the
application that prompted it, which is why it is written up here in the notation of that paper.
Three things come out of it.

First, an exact identity.
Because the integrand $\abs{Z}^p \abs{1 - Z}^q$ is the squared modulus of the holomorphic function
$f = Z^{p/2} (1 - Z)^{q/2}$, Stokes' theorem turns \eqref{jb:eq:J-def} into a contour integral, and
$J(p, q)$ becomes the area swept by the image of $R$ under an antiderivative $G$ of $f$ -- which is
an incomplete Beta function, hence a Gauss ${}_2F_1$.
This is \cref{jb:sec:stokes}, and it is the whole content of the note in one line.

Second, an arc-by-arc analysis (\cref{jb:sec:arcs,jb:sec:obstruction}).
The real segment contributes exactly zero for all $(p, q)$.
Each of the other three arcs has a straight image under $G$, and then closes in endpoint values of
$G$, on one locus of the $(p, q)$ plane: the vertical segment $\Re Z = \tfrac12$ on $p = q$, the
unit circle on $2p + q = -4$ (where it contributes exactly zero), the ellipse only at the point
$p = q = -1$.
The two curved arcs -- the unit circle and the ellipse -- therefore do not collapse for generic
$(p, q)$, and it is worth knowing why, because the reason is not the one a first look suggests.
It is \emph{not} the absence of a symmetry: each arc carries an algebraic Schwarz function, so
complex conjugation of $G$ along the arc is the restriction of a holomorphic function and every arc
contribution is an honest contour integral of a holomorphic integrand.
Nor is it the awkward endpoints $\e^{\I \pi/3}$, $\tfrac12 + \I b$, which are not special points of
the hypergeometric function.
What blocks the evaluation is that the resulting integrand carries \emph{two} algebraic powers
against the ${}_2F_1$, with exponents unrelated to its parameters; the antiderivative of such a
product is a double series, and it collapses to a single hypergeometric series when Euler's
transformation can absorb one of the powers -- a sufficient condition, met on a few lines, and one
that closes the circle arc only.
On the lines $p = 0$ and $q = 0$, where $G$ is elementary, the vertex integral does close, and we
give $J(0, q)$ and $J(p, 0)$ with every coefficient, in incomplete Beta functions and Appell $F_1$,
checked numerically to $10^{-15}$; at the single point $p = q = -1$ all four arcs close at once and
$J$ is a combination of ${}_2F_1$ endpoint values and one ${}_3F_2$.
Everywhere else the honest statement is: an exact one-dimensional integral of a ${}_2F_1$, no
closed form found by this route.
This is a scoped negative statement about one route, not a proof that no closed form exists.

Third, a structural observation about the vertex region itself (\cref{jb:sec:confocal}), independent of
the boundary reduction.
The semi-axes of the $s$-channel ellipse satisfy $a^2 - b^2 = \tfrac14$, which is exactly the
statement that the foci sit at the punctures $0$ and $1$.
Consequently $\partial \mc V_{0,4}$ is a level set of $\abs{Z} + \abs{1 - Z}$, and in the variables
$u = \abs{Z} + \abs{1 - Z}$, $v = \abs{1 - Z} - \abs{Z}$ the fundamental domain of the vertex region
is a trapezoid with an elementary algebraic weight.
Three of its four sides are coordinate lines; the exception is the unit circle.
This isolates the unit-circle constraint $\abs{Z} \ge 1$ -- not the vertex boundary -- as the
single feature of the region that obstructs separation of variables.

\Cref{jb:sec:numerics,jb:sec:use} are practical: the boundary identity evaluated by cumulative Chebyshev
integration needs no hypergeometric function at all, costs a few thousand elementary evaluations,
and reproduces the two-dimensional quadrature to its own accuracy at exponents as large as
$\abs{p} + \abs{q} \approx 170$.
That is the form in which we recommend the result be used.



\subsection{The region and the integral}
\label{jb:sec:setup}


\subparagraph{The region.}
We work inside the $S_4$ fundamental domain of the vertex region, defined in
\cite{Erbin:2023:SL2CQuarticVertex} by
\begin{equation}
  \label{jb:eq:region}
  R
    = \Acom{
      \Im Z \ge 0,
      \quad
      \abs{Z} \ge 1,
      \quad
      \Re Z \le \tfrac12,
      \quad
      \frac{(x - \tfrac12)^2}{a^2} + \frac{y^2}{b^2} \le 1
    },
\end{equation}
the last condition being the interior of the $s$-channel boundary curve, an ellipse centered at
$z_c = \tfrac12$ with semi-axes
\begin{equation}
  \label{jb:eq:axes}
  a
    = \beta^2 + \frac{1}{16 \beta^2},
  \qquad
  b
    = \beta^2 - \frac{1}{16 \beta^2},
  \qquad
  a^2 - b^2
    = \frac14 .
\end{equation}
Here $\beta$ is the vertex parameter of \cite{Erbin:2023:SL2CQuarticVertex}, related to the stub
length; we take $\beta \ge \beta_{04} := \tfrac12 + \tfrac{1}{\sqrt2} \approx 1.207$ throughout,
the range on which the four-arc description of $\partial R$ below holds: it is exactly the
condition $a \ge \tfrac32$ that the ellipse meets the real axis at or to the left of $-1$, and it
implies $b \ge \tfrac{\sqrt3}{2}$, so that the vertical side reaches above $\e^{\I \pi/3}$.
Every numerical check uses the test values $\beta = \tfrac32$ and $\beta = 2$.
The relation $a^2 - b^2 = \tfrac14$ is not an extra input: it is the value $c = \tfrac12$ of the
focal distance already recorded with the ellipse parameters in
\cite{Erbin:2023:SL2CQuarticVertex}.
\Cref{jb:sec:confocal} is nothing but the consequence of reading it as a statement about the foci.

The boundary $\partial R$ consists of the four curves $C_1$--$C_4$ of
\cite{Erbin:2023:SL2CQuarticVertex}.
We orient it positively, interior on the left, and start at $Z = -1$:
\begin{enumerate}
  \setlength{\itemsep}{0.3ex}

  \item the unit circle $C_3$ from $-1$ to $\e^{\I \pi/3}$ (traversed clockwise on the circle,
    $\theta : \pi \to \pi/3$);

  \item the vertical segment $C_2$, $\Re Z = \tfrac12$, from $\e^{\I \pi/3}$ up to
    $\tfrac12 + \I b$;

  \item the elliptic quarter $C_4$ from $\tfrac12 + \I b$ to $\tfrac12 - a$
    ($\theta : \tfrac{\pi}{2} \to \pi$ in the eccentric-anomaly parametrization);

  \item the real segment $C_1$ from $\tfrac12 - a$ rightwards to $-1$.
\end{enumerate}
$R$ is compact, simply connected, contained in the closed upper half plane, and avoids the
punctures $Z = 0$ and $Z = 1$; on it $1 \le \abs{Z} \le a$ and $1 \le \abs{1 - Z} \le a + \tfrac12$,
so $J(p, q)$ is an entire function of $(p, q)$.
Note that $R$ touches the negative real axis along $C_1$, which is the only delicate point in what
follows.

\subparagraph{Branches.}
The integrand of \eqref{jb:eq:J-def} is the squared modulus of
\begin{equation}
  \label{jb:eq:f-def}
  f(Z)
    := Z^{p/2} \, (1 - Z)^{q/2},
\end{equation}
and everything below depends on choosing the two powers so that $f$ is single-valued and continuous
on all of $R$, including on the negative axis.
On the closed upper half plane minus $\{0, 1\}$ we set
\begin{subequations}
  \label{jb:eq:branches}
\begin{align}
  Z^{p/2}
    &
    := \abs{Z}^{p/2} \, \e^{\I p \theta_0 / 2},
  \qquad
  \theta_0
    = \arg Z
    \in [0, \pi],
  \\
  (1 - Z)^{q/2}
    &
    := \abs{1 - Z}^{q/2} \, \e^{\I q \theta_1 / 2},
  \qquad
  \theta_1
    = \arg(1 - Z)
    \in [-\pi, 0].
\end{align}
\end{subequations}
Both are the principal branches, with cuts on $(-\infty, 0]$ and $[1, \infty)$ respectively; the
first is taken as its boundary value \emph{from above} on the negative axis.
So defined, $f$ is holomorphic in the open upper half plane, continuous up to its boundary away from
$0$ and $1$, and
\begin{equation}
  \label{jb:eq:modulus}
  f \bar f
    = \abs{Z}^{p} \abs{1 - Z}^{q},
\end{equation}
the phases canceling identically.
The reason to belabor this is that $R$ has a boundary arc \emph{on} the cut of $Z^{p/2}$: the
boundary value from above must be used consistently in $f$ and in its antiderivative, or the
reduction below is off by a phase on that arc.

\subparagraph{Exponents.}
We take $p, q$ real, and exclude $p \in \{-2, -4, \dots\}$, where the normalization of the
antiderivative used below has a pole.
The identity extends to those values by continuity, since $J$ is entire; see the remark after
\eqref{jb:eq:G-def}.


\subsection{The tile integral as a swept area}
\label{jb:sec:stokes}


This section proves the one identity everything else rests on: the two-dimensional integral
\eqref{jb:eq:J-def} equals a contour integral along $\partial R$, and has a geometric reading as an
area in the image plane of an antiderivative of $f$.

\subparagraph{Setup.}
Assume \eqref{jb:eq:region}--\eqref{jb:eq:branches}.
Write
\begin{equation}
  \label{jb:eq:alpha-gamma}
  \alpha
    := \frac{p}{2} + 1,
  \qquad
  \gamma
    := \frac{q}{2} + 1,
\end{equation}
and let $G$ be any holomorphic antiderivative of $f$ on the open upper half plane $\mathbb{H}$,
$G' = f$.

\subparagraph{Derivation.}
Consider the $1$-form $\omega = f \, \overline{G} \, \dd Z$ on $\mathbb{H}$.
It is not closed, and its failure to be closed is exactly the integrand we want: since $f$ is
holomorphic, only the $\bar Z$-derivative of $\overline{G}$ survives,
\begin{equation}
  \label{jb:eq:domega}
  \dd \omega
    = \pd_{\bar Z} \bigl( f \overline{G} \bigr) \, \dd \bar Z \wedge \dd Z
    = f \, \overline{G'} \, \dd \bar Z \wedge \dd Z
    = 2 \I \, \abs{f}^2 \, \dd x \wedge \dd y,
\end{equation}
where we used $\dd \bar Z \wedge \dd Z = 2 \I \, \dd x \wedge \dd y$ and, in the last step,
\eqref{jb:eq:modulus}.
The region $R$ is compact and simply connected with piecewise-smooth boundary, and $\omega$ is
continuous up to that boundary, so Stokes' theorem applies and gives the identity we are after:
\begin{equation}
  \label{jb:eq:stokes}
  J(p, q)
    = \frac{1}{2 \I} \oint_{\partial R} f \, \overline{G} \, \dd Z
    = \frac12 \, \Im \oint_{\partial R} \overline{G} \, \dd G .
\end{equation}
The second form follows from $\dd \abs{G}^2 = \overline{G} \, \dd G + G \, \overline{\dd G}$: the
combination $\overline{G} \, \dd G$ splits into $\tfrac12 \dd \abs{G}^2$ plus $\I$ times its own
imaginary part, and the exact piece integrates to zero around a closed curve.

The second form of \eqref{jb:eq:stokes} is the useful one, because it is manifestly real and has a
geometric meaning: $\tfrac12 \Im \int \overline{G} \, \dd G$ over a curve is the area swept from the
origin by the moving radius vector of the image curve.
Hence
\begin{quote}
  \emph{$J(p, q)$ is the area of the image $G(R)$, counted with multiplicity.}
\end{quote}
In particular the contribution of each boundary arc is the area it sweeps in the $G$-plane,
\begin{equation}
  \label{jb:eq:per-arc}
  J
    = \sum_k J_k,
  \qquad
  J_k
    = \frac12 \int_{t_0^{(k)}}^{t_1^{(k)}}
      \Im \bigl( \overline{G} \, G'(Z) \, Z'(t) \bigr) \, \dd t,
\end{equation}
for any parametrization $t \mapsto Z(t)$ of arc $k$.
The differential $\dd Z$ is complex on three of the four arcs, so it must stay inside the imaginary
part.
It is the swept area \eqref{jb:eq:per-arc}, and not the naive per-arc piece of the first form of
\eqref{jb:eq:stokes}, that is the meaningful quantity arc by arc: splitting
$\tfrac{1}{2\I} \int f \overline{G} \, \dd Z$ leaves on each arc a spurious \emph{imaginary}
contribution $\tfrac{1}{4\I} [\abs{G}^2]$, real bracket times imaginary prefactor, which telescopes
to zero only around the full loop.
The additive constant in $G$ drops out of \eqref{jb:eq:stokes} by Cauchy's theorem
($\oint_{\partial R} f \, \dd Z = 0$); we fix it by
\begin{equation}
  \label{jb:eq:base-point}
  G(-1)
    = 0 .
\end{equation}

\subparagraph{The antiderivative.}
An antiderivative of $Z^{\alpha - 1} (1 - Z)^{\gamma - 1}$ with the branches \eqref{jb:eq:branches} is
the incomplete Beta function,
\begin{equation}
  \label{jb:eq:incomplete-beta}
  B_Z(\alpha, \gamma)
    = \frac{Z^{\alpha}}{\alpha} \, {}_2F_1(\alpha, 1 - \gamma; \alpha + 1; Z),
  \qquad
  \frac{\dd}{\dd Z} B_Z(\alpha, \gamma)
    = Z^{\alpha - 1} (1 - Z)^{\gamma - 1},
\end{equation}
so that, with the normalization \eqref{jb:eq:base-point},
\begin{equation}
  \label{jb:eq:G-def}
  G(Z)
    = B_Z(\alpha, \gamma) - B_{-1}(\alpha, \gamma)
    = \int_{-1}^{Z} f,
\end{equation}
the integral being along any path in the closed upper half plane avoiding $0$ and $1$.
The hypergeometric representation \eqref{jb:eq:incomplete-beta} is standard \cite{Nist:DLMF}; the point
to check is that it is compatible with our branch choice.
The ${}_2F_1$ continued outside the unit disk has its cut on $[1, \infty)$, the same as
$(1 - Z)^{\gamma - 1}$, and on $(-\infty, 0]$ we take its boundary value from above, matching
\eqref{jb:eq:branches}; with that reading \eqref{jb:eq:incomplete-beta} is an identity of analytic
functions wherever $\alpha \notin \{0, -1, -2, \dots\}$.

Two remarks on the exceptional exponents.
For $p \le -2$ the right-hand side of \eqref{jb:eq:incomplete-beta} is no longer the convergent integral
$\int_0^Z w^{\alpha - 1} (1 - w)^{\gamma - 1} \dd w$, but it is still an antiderivative of $f$, which
is all \eqref{jb:eq:stokes} needs.
At $\alpha \in \{0, -1, -2, \dots\}$ the normalization $Z^\alpha / \alpha$ has a pole; the difference
\eqref{jb:eq:G-def} has a finite limit there, and $J$ is entire, so nothing is lost.

\subparagraph{Checks.}
Three independent numerical checks confirm this section.
The identity \eqref{jb:eq:incomplete-beta} was verified by differentiating $G$ numerically at $30$-digit
precision at eight interior points with $\abs{p}, \abs{q} < 30$ (agreement to $\sim 10^{-26}$ or
better), and $G$ from \eqref{jb:eq:G-def} was compared to the direct path integral $\int_{-1}^Z f$
(agreement to $5 \times 10^{-30}$).
Around the closed boundary, $\Re \oint f \overline{G} \, \dd Z$ vanishes and $G$ returns to zero to
$7 \times 10^{-15}$, which tests the orientation, the arc parametrizations, and the branch
bookkeeping at once.
The global accuracy of \eqref{jb:eq:stokes} against two-dimensional quadrature is reported in
\cref{jb:sec:numerics}.


\subsection{The four arcs}
\label{jb:sec:arcs}


We now evaluate the four contributions \eqref{jb:eq:per-arc} one at a time.
The organizing idea is elementary.
The image of an arc under $G$ is traced by $\dd G = f \, \dd Z$, so it is a straight segment if
and only if the \emph{direction} of $\dd G$, that is $\arg(f \, \dd Z)$, is constant along the arc.
On the two straight arcs $\dd Z$ has a fixed phase and the criterion reduces to a constant
$\arg f$; on the two curved arcs the phase of $\dd Z$ turns with the arc and must be counted.
A straight image sweeps an area that is elementary to write down: \emph{no} area when its line
passes through the base point $G(-1) = 0$, and otherwise the triangle
$\tfrac12 \Im \bigl( \overline{G(Z_0)} \, G(Z_1) \bigr)$ spanned by the origin and the two
endpoint values, $Z_0$ and $Z_1$ being the start and the end of the arc in the orientation of
$\partial R$.
Call this the \emph{constant-phase mechanism}.
It closes the real segment for every $(p, q)$, and each of the other three arcs on one locus in the
$(p, q)$ plane -- the vertical segment on the line $p = q$, the unit circle on the line
$2 p + q = -4$, the ellipse at the single point $p = q = -1$.
At generic $(p, q)$ it therefore closes one arc, and the failure on the other three is instructive.


\subsubsection{The real segment: exactly zero}
\label{jb:sec:arcs:real}

On $C_1$ we have $Z = x \in [\tfrac12 - a, -1]$, hence $\theta_0 = \pi$ and $\theta_1 = 0$ in
\eqref{jb:eq:branches}, and $\dd Z = \dd x$ is real.
The phase of $f \, \dd Z$ is therefore that of $f$, and it is constant along the whole segment:
\begin{equation}
  \label{jb:eq:f-real}
  f(x)
    = \e^{\I \pi p / 2} \, \phi(x),
  \qquad
  \phi(x)
    = \abs{x}^{p/2} (1 - x)^{q/2}
    > 0 .
\end{equation}
Integrating from the base point, $G(x) = - \e^{\I \pi p/2} \Phi(x)$ with
$\Phi(x) = \int_x^{-1} \phi$ real and positive.
The image of $C_1$ is thus the radial ray $\{ - \e^{\I \pi p/2} \Phi \}$ emanating from
$G(-1) = 0$, on which $\overline{G} \, G'$ is real pointwise, and
\begin{equation}
  \label{jb:eq:J-real}
  J_{\text{real}}
    = 0 .
\end{equation}
This is the arc that sits on the cut of $Z^{p/2}$; it is precisely the branch prescription
\eqref{jb:eq:branches} that makes its image a straight ray, and the contribution vanishes for all
$(p, q)$.
Note that the vanishing is an artifact of the base point, not of the arc: with a different
normalization the segment would contribute
$\tfrac12 \Im \bigl( \overline{G(-1)} \, \e^{\I \pi p/2} \bigr) \Phi(\tfrac12 - a)$ with
$\Phi(\tfrac12 - a) = \int_1^{a - 1/2} t^{p/2} (1 + t)^{q/2} \, \dd t$ a real incomplete Beta
function -- closed either way, but the choice \eqref{jb:eq:base-point} makes it free.
Numerically the swept area of this arc relative to $J$ is $2 \times 10^{-34}$.


\subsubsection{The vertical segment: a straight image exactly when $p = q$}
\label{jb:sec:arcs:vertical}

On $C_2$, parametrize $Z = \tfrac12 + \I y$ with $y \in [\tfrac{\sqrt3}{2}, b]$.
This is the line equidistant from the two punctures, so $1 - Z = \bar Z$ and
\begin{equation}
  \label{jb:eq:f-vert}
  f
    = \abs{Z}^{(p + q)/2} \, \e^{\I (p - q) \theta_0 / 2},
  \qquad
  \theta_0
    = \arctan 2y
    \in \Bigl[ \frac{\pi}{3}, \frac{\pi}{2} \Bigr).
\end{equation}
Here $\dd Z = \I \, \dd y$, so $\arg(f \, \dd Z) = \tfrac{(p - q) \theta_0}{2} + \tfrac{\pi}{2}$
is constant along the segment if and only if $p = q$, and the mechanism of
\cref{jb:sec:arcs:real} then applies in rotated form.
For $p = q$ the function $f = \abs{Z}^p$ is real, so $G$ moves vertically in the image plane:
$G(Z) = G_1 + \I \Psi(y)$ with $G_1 = G(\e^{\I \pi/3})$ the value at the lower endpoint and
$\Psi$ real.
Then $\overline{G} \, \dd G = (\overline{G_1} - \I \Psi) \, \I \abs{Z}^p \, \dd y$, whose imaginary
part is $\Re(G_1) \abs{Z}^p \dd y$, and
\begin{equation}
  \label{jb:eq:J-vert}
  J_{\text{vert}} \big|_{p = q}
    = \frac12 \, \Re (G_1) \, \Psi(b),
  \qquad
  \Psi(b)
    = \int_{\sqrt3/2}^{b} \Bigl( \frac14 + y^2 \Bigr)^{p/2} \dd y,
\end{equation}
the remaining integral being a Gauss ${}_2F_1$.
For $p \ne q$ the phase $\e^{\I (p - q) \theta_0 / 2}$ varies along the segment, the image is a
curved arc, and \eqref{jb:eq:J-vert} is simply wrong.
The segment then still closes, as an incomplete Beta function, on the three lines where $G$
itself is elementary -- $p = 0$, $q = 0$ and $p + q = -4$ (\cref{jb:sec:obstruction}) -- but off
those lines no reduction of
$\int \Im \bigl( \overline{G} \, G'(Z) \, Z'(t) \bigr) \dd t$ below a one-dimensional integral of a
${}_2F_1$ was found on this arc.
Both statements about \eqref{jb:eq:J-vert} are checked numerically: it holds to $2 \times 10^{-17}$ at
$p = q$, and applying it at $p \ne q$ produces a relative deviation of at least $0.47$.


\subsubsection{The curved arcs: constant phase on a line, and at a point}
\label{jb:sec:arcs:curved}

On the unit circle $C_3$, write $Z = \e^{\I \theta}$ with $\theta \in [\pi/3, \pi]$.
Using $1 - \e^{\I \theta} = 2 \sin \tfrac{\theta}{2} \, \e^{\I (\theta - \pi)/2}$, which gives
$\theta_1 = \tfrac{\theta - \pi}{2} \in [-\tfrac{\pi}{3}, 0]$, consistent with
\eqref{jb:eq:branches}, we get
\begin{equation}
  \label{jb:eq:f-circle}
  f
    = \Bigl( 2 \sin \frac{\theta}{2} \Bigr)^{q/2} \,
      \e^{\I [ 2 p \theta + q (\theta - \pi) ] / 4},
  \qquad
  \dd Z
    = \I \e^{\I \theta} \, \dd \theta ,
\end{equation}
the factor $\tfrac{q}{4}$ being the $\tfrac{q}{2}$ of $(1 - Z)^{q/2}$ applied to
$\theta_1 = \tfrac{\theta - \pi}{2}$.
The phase of $f$ is linear in $\theta$ with slope $\tfrac{p}{2} + \tfrac{q}{4}$, and the factor
$\dd Z$ adds a further $\e^{\I \theta}$, so
\begin{equation}
  \label{jb:eq:phase-circle}
  \arg (f \, \dd Z)
    = \theta \Bigl( \frac{p}{2} + \frac{q}{4} + 1 \Bigr)
      + \frac{\pi}{2}
      - \frac{\pi q}{4} .
\end{equation}
This is constant along the arc if and only if $2 p + q = -4$.
On that line the image of $C_3$ is a straight segment starting at $G(-1) = 0$, so its line passes
through the origin and
\begin{equation}
  \label{jb:eq:J-circle-line}
  J_{\text{circle}} \big|_{2 p + q = -4}
    = 0
\end{equation}
exactly, for every $\beta$.
Note that the phase is \emph{not} constant at $(p, q) = (0, 0)$, where the slope in
\eqref{jb:eq:phase-circle} is $1$: the factor $\dd Z$ alone turns the image.
Numerically, the spread of the unit phase $f \, \dd Z / \abs{f \, \dd Z}$ along the arc is below
$2 \times 10^{-16}$ at seven points on the line $2 p + q = -4$, with $\abs{J_{\text{circle}} / J}$
below $4 \times 10^{-17}$; off the line, at $(-1, -1)$ and $(0, 0)$ for instance, the spread is
$0.47$ and $1.62$ and the circle contributes $-0.13 \, J$ and $-0.23 \, J$.

On the elliptic quarter $C_4$ the natural parametrization is the one adapted to the foci.
Write $Z = \cosh^2 \tfrac{\zeta}{2}$ with $\zeta = \xi_0 + \I \eta$ and $\xi_0 = 2 \ln 2 \beta$
fixed: this is the ellipse \eqref{jb:eq:axes} with foci at $0$ and $1$, and the quarter $C_4$
corresponds to $\tfrac{\eta}{2} \in [\tfrac{\pi}{4}, \tfrac{\pi}{2}]$.
Then $1 - Z = - \sinh^2 \tfrac{\zeta}{2}$ and
$\dd Z = \I \sinh \tfrac{\zeta}{2} \cosh \tfrac{\zeta}{2} \, \dd \eta$.
On this range of $\eta$ both $\cosh \tfrac{\zeta}{2}$ and $\sinh \tfrac{\zeta}{2}$ lie in the
closed first quadrant, so the principal branches \eqref{jb:eq:branches} give
$Z^{p/2} = \cosh^p \tfrac{\zeta}{2}$ and $(1 - Z)^{q/2} = \e^{- \I \pi q/2} \sinh^q \tfrac{\zeta}{2}$,
and
\begin{equation}
  \label{jb:eq:phase-ellipse}
  f \, \dd Z
    = \I \e^{- \I \pi q / 2} \, \cosh^{p + 1} \frac{\zeta}{2} \, \sinh^{q + 1} \frac{\zeta}{2} \,
      \dd \eta,
  \qquad
  \arg (f \, \dd Z)
    = (p + 1) \arctan \Bigl( \tanh \frac{\xi_0}{2} \tan \frac{\eta}{2} \Bigr)
      + (q + 1) \arctan \Bigl( \coth \frac{\xi_0}{2} \tan \frac{\eta}{2} \Bigr)
      + \text{const}.
\end{equation}
The $\eta$-derivative of this phase is a rational function of $T = \tan \tfrac{\eta}{2}$, which
vanishes on the interval $T \in [1, \infty)$ only if it vanishes identically; its values at $T = 0$
and $T \to \infty$ give the homogeneous system
$(p + 1) \tanh x + (q + 1) \coth x = 0$, $(p + 1) \coth x + (q + 1) \tanh x = 0$ with
$x = \tfrac{\xi_0}{2} > 0$, whose determinant $\tanh^2 x - \coth^2 x$ is negative.
So the phase is constant on the ellipse if and only if
\begin{equation}
  \label{jb:eq:ellipse-point}
  p
    = q
    = -1,
\end{equation}
a single point, where $f \, \dd Z = - \dd \eta$ is real.
There the image is the straight segment from $G(\tfrac12 + \I b)$ to $G(\tfrac12 - a)$, which does
not pass through the origin, and the arc contributes the triangle
\begin{equation}
  \label{jb:eq:J-ellipse-point}
  J_{\text{ellipse}} \big|_{p = q = -1}
    = \frac12 \, \Im \Bigl( \overline{G(\tfrac12 + \I b)} \, G(\tfrac12 - a) \Bigr),
\end{equation}
elementary in two endpoint values of $G$, each a ${}_2F_1$.
Numerically, at $(-1, -1)$ and $\beta = \tfrac32$ the phase spread is $1.5 \times 10^{-16}$ and
\eqref{jb:eq:J-ellipse-point} agrees with the swept area to better than $10^{-15}$; at the nearby
points $(-1, -1.3)$ and $(-0.7, -1)$ the spread is $0.18$ and $0.24$ and the formula is off by
$14 \%$ and $15 \%$.

The count is therefore as follows.
At generic $(p, q)$ the constant-phase mechanism closes one arc, the real segment.
On each of the loci $p = q$, $2 p + q = -4$ and $p = q = -1$ it closes a second one; at the
intersections $(-\tfrac43, -\tfrac43)$ and $(-1, -1)$ it closes three; it never closes all four,
since $p = q = -1$ and $2 p + q = -4$ are incompatible.
None of these loci closes $J$ by itself.
The one point where $J$ does close entirely, $(-1, -1)$, needs one more ingredient for the circle
arc, given in \cref{jb:sec:obstruction:powers}.

This is where one would look for a second mechanism, and there is one -- reflection.
It works, and it still does not produce a closed form; the next section explains both halves of that
statement.


\subsection{Reflection, and what actually obstructs a closed form}
\label{jb:sec:obstruction}


The obstruction to closing $C_3$ and $C_4$ is easy to misdiagnose.
The natural guess is that $\overline{G}$, being anti-holomorphic, cannot be handled by contour
methods, so one needs a symmetry of the region that $\mathbb{C} \setminus \{0, 1\}$ does not have.
That guess is wrong.
Every arc of $\partial R$ admits a Schwarz function, complex conjugation along each arc \emph{is}
the restriction of a holomorphic map, and after using it each arc contribution becomes an ordinary
contour integral of a holomorphic integrand.
The obstruction lies one step further on, in the integrand that results.


\subsubsection{Schwarz functions of the four arcs}
\label{jb:sec:obstruction:schwarz}

The \emph{Schwarz function} of an analytic arc is the analytic function $S$ that agrees with
$Z \mapsto \bar Z$ on that arc and nowhere else \cite{Davis:1974:SchwarzFunction}.
It exists and is algebraic for lines, circles and conics.
For our four arcs,
\begin{subequations}
  \label{jb:eq:schwarz}
\begin{gather}
  S(Z)
    = \frac{1}{Z} \ (C_3),
  \qquad
  S(Z)
    = 1 - Z \ (C_2),
  \qquad
  S(Z)
    = Z \ (C_1),
  \\
  S(Z)
    = \frac12 + 4 (a^2 + b^2) \Bigl( Z - \frac12 \Bigr) + 8 a b \sqrt{Z (Z - 1)} \ (C_4).
\end{gather}
\end{subequations}
The last is two-valued; the sign written is the sheet selected by the principal square root along
the open quarter from $\tfrac12 - a$ to $\tfrac12 + \I b$.
At the endpoint $\tfrac12 + \I b$ the argument $Z(Z - 1) = -a^2$ lands on the cut of the principal
root and the value must be taken as the limit along the arc; this is a single point and plays no
role.
It is worth pausing on the square root in \eqref{jb:eq:schwarz}: its branch points sit at $Z = 0$ and
$Z = 1$, the two punctures, and they do so \emph{because} $a^2 - b^2 = \tfrac14$.
The same identity reappears, in a different guise, in \cref{jb:sec:confocal}.

Now use the reflection property of the incomplete Beta function: for real $p, q$ one has
$\overline{B_{\bar w}(\alpha, \gamma)} = B_w(\alpha, \gamma)$ on the cut plane.
Combining this with $\bar Z = S(Z)$ on the arc gives
\begin{equation}
  \label{jb:eq:reflection}
  \overline{G(Z)}
    = G \bigl( S(Z) \bigr) + \text{const}
  \qquad \text{on each curved arc,}
\end{equation}
the right-hand side being holomorphic in $Z$ in a neighborhood of the arc -- on the two-sheeted
algebraic curve of $S$, in the elliptic case.
Every arc contribution to \eqref{jb:eq:stokes} is therefore a contour integral of a holomorphic
integrand.
Identity \eqref{jb:eq:reflection} has been verified on all three curved arcs: the difference
$\overline{G(Z)} - G(S(Z))$ is constant along each arc to a relative spread of $6 \times 10^{-29}$,
and feeding an arc the Schwarz function of a \emph{different} arc destroys the constancy, as it
must.

One arc is an exception, and it is the one on the cut.
On the real segment $C_1$ the Schwarz function is $S(Z) = Z$, and \eqref{jb:eq:reflection} holds only
if $G(S(Z))$ is read as the boundary value of $B_w$ from \emph{below} the cut,
$B_x^{\text{below}} = \overline{B_x^{\text{above}}}$.
Our branch choice \eqref{jb:eq:branches} fixes $G$ by its value from above, and the two differ there by
the phase $\e^{-\I \pi p}$, which is not a constant of the required kind.
Nothing rests on this: $C_1$ is closed exactly by \eqref{jb:eq:J-real}.


\subsubsection{The reduced circle integral}
\label{jb:sec:obstruction:circle}

Apply \eqref{jb:eq:reflection} to the circle, substituting $w = S(Z) = 1/Z$ in
$f \, \overline{G} \, \dd Z$.
Two simplifications occur.
The powers recombine: with $Z = \e^{\I \theta}$, $\theta \in (\tfrac{\pi}{3}, \pi)$, one has
$\arg w = -\theta \in (-\pi, -\tfrac{\pi}{3})$, still principal, and both
$Z^{p/2} Z^{-\alpha} = Z^{-1}$ and $w^{-q/2} (w - 1)^{q/2} = (1 - Z)^{q/2}$ hold with all branches
principal.
The Jacobian is $\dd Z = - \dd w / w^2$, whence the overall sign.
The result is
\begin{equation}
  \label{jb:eq:circle-reduced}
  \frac{1}{2 \I} \int_{C_3} f \, \overline{G} \, \dd Z
    = - \frac{1}{2 \I \alpha} \int_{-1}^{\e^{- \I \pi/3}}
      w^{- q/2 - 1} (w - 1)^{q/2} \,
      {}_2F_1(\alpha, 1 - \gamma; \alpha + 1; w) \, \dd w
      - \frac{\overline{B_{-1}(\alpha, \gamma)}}{2 \I}
      \bigl[ G \bigr]_{-1}^{\e^{\I \pi/3}},
\end{equation}
all branches principal, the second term a pure endpoint value.
The left-hand side is the circle piece of the \emph{first} form of \eqref{jb:eq:stokes}, and, as
noted after \eqref{jb:eq:per-arc}, it is not the swept area of the arc: splitting
$\overline{G} \, \dd G = \tfrac12 \dd \abs{G}^2 + \I \Im(\overline{G} \, \dd G)$ on the arc gives
\begin{equation}
  \label{jb:eq:circle-first-form}
  \frac{1}{2 \I} \int_{C_3} f \, \overline{G} \, \dd Z
    = J_{\text{circle}}
      + \frac{1}{4 \I} \bigl[ \abs{G}^2 \bigr]_{-1}^{\e^{\I \pi/3}}
    = J_{\text{circle}}
      - \frac{\I}{4} \, \abs{G(\e^{\I \pi/3})}^2,
\end{equation}
since $G(-1) = 0$.
The circle's swept area is thus the real part of \eqref{jb:eq:circle-reduced}; the imaginary
remainder telescopes around the loop and nothing in $J$ depends on it.
This is checked at two points, $(-3.4, 2.1)$ and $(1.7, -5.3)$ at $\beta = \tfrac32$: the
right-hand side of \eqref{jb:eq:circle-reduced} by quadrature along $w = \e^{-\I\theta}$ differs from
the swept area by a purely imaginary number (real part of the difference below
$2 \times 10^{-16}$) equal to $- \tfrac{\I}{4} \abs{G(\e^{\I \pi/3})}^2$ to the digits shown,
$-3.12795 \, \I$ and $-0.103872 \, \I$.
The same manipulation on the vertical segment, integrating by parts, gives
$J_{\text{vert}} = \tfrac{1}{2 \I} ( [G \tilde G] - \int G \, \tilde f \, \dd Z + c \, [G] )$, where
$\tilde G(Z) = B_{1 - Z}(\alpha, \gamma)$ is the antiderivative with $p$ and $q$ exchanged,
$\tilde f = \tilde G' = - Z^{q/2} (1 - Z)^{p/2}$, and $c$ is the constant in \eqref{jb:eq:reflection}.
The elliptic arc works the same way through the two-valued $S$ of \eqref{jb:eq:schwarz}.

So the reduction runs to completion: the curved arcs are contour integrals of
(algebraic powers) $\times$ ${}_2F_1$.
The question is whether such an integral has an antiderivative in the classes we are willing to
accept -- finite combinations of Gauss ${}_2F_1$ and Appell $F_1$, $F_2$.


\subsubsection{What obstructs: a pair of powers, not the endpoints}
\label{jb:sec:obstruction:powers}

It is tempting to blame the endpoints of \eqref{jb:eq:circle-reduced}, which are $-1$ and
$\e^{-\I\pi/3}$ rather than the hypergeometric function's natural $0$, $1$, $\infty$.
They are not the problem.
A \emph{single} power against a ${}_2F_1$ integrates term by term to a ${}_3F_2$ at an arbitrary
endpoint,
\begin{equation}
  \label{jb:eq:one-power}
  \int_0^z w^{\lambda} \, {}_2F_1(A, B; C; w) \, \dd w
    = \frac{z^{\lambda + 1}}{\lambda + 1} \,
      {}_3F_2(A, B, \lambda + 1; C, \lambda + 2; z),
\end{equation}
which is a closed form at any $z$.
The obstruction in \eqref{jb:eq:circle-reduced} is the \emph{second} power, $(w - 1)^{q/2}$.
With both powers present, expanding the hypergeometric series and the binomial and integrating term
by term gives
\begin{equation}
  \label{jb:eq:two-powers}
  \int_0^z w^{\lambda} (1 - w)^{\mu} \, {}_2F_1(A, B; C; w) \, \dd w
    = \frac{z^{\lambda + 1}}{\lambda + 1} \sum_{m, n \ge 0}
      \frac{(\lambda + 1)_{m + n}}{(\lambda + 2)_{m + n}} \,
      \frac{(- \mu)_m \, (A)_n (B)_n}{(C)_n \, m! \, n!} \, z^{m + n},
\end{equation}
a double series in the \emph{single} variable $z$; at $\mu = 0$ only $m = 0$ survives and it
reduces to \eqref{jb:eq:one-power}.
This is a Kampé de Fériet function, not an Appell function
\cite{Srivastava:1985:MultipleGaussianHypergeometric}: the Pochhammer factor
$(A)_n (B)_n / (C)_n$ carries one summation index only, so the series does not factorize into the
Appell pattern, and the two indices are locked together by $(\lambda + 1)_{m+n}/(\lambda+2)_{m+n}$.
It collapses to a single hypergeometric series when Euler's transformation absorbs the extra
power: rearranging
${}_2F_1(A, B; C; w) = (1 - w)^{C - A - B} \, {}_2F_1(C - A, C - B; C; w)$ gives
\begin{equation}
  \label{jb:eq:euler-condition}
  (1 - w)^{\mu} \, {}_2F_1(A, B; C; w)
    = {}_2F_1(C - A, C - B; C; w)
  \qquad \text{for} \quad
  \mu
    = A + B - C,
\end{equation}
after which \eqref{jb:eq:one-power} applies.
(With the opposite sign, $\mu = C - A - B$, the product is
$(1 - w)^{2 \mu} {}_2F_1(C - A, C - B; C; w)$ and nothing is absorbed.)
Absorption is a sufficient condition for a single series, not a necessary one: we know of no
necessary condition, and the degenerate case below closes the arc without absorbing anything.

Apply \eqref{jb:eq:euler-condition} to the integrand
$I(w) \, \dd w = w^{-\gamma} (w - 1)^{\gamma - 1} {}_2F_1(\alpha, 1 - \gamma; \alpha + 1; w) \, \dd w$
of \eqref{jb:eq:circle-reduced}.
There are three ways to do so, one per representative of the ${}_2F_1$.
Under the Pfaff substitution $x = w / (w - 1)$, that is $w = - x / (1 - x)$, $1 - w = (1 - x)^{-1}$,
$\dd w = - \dd x / (1 - x)^2$, the arc $w = \e^{-\I\theta}$, $\theta \in (\tfrac{\pi}{3}, \pi)$,
becomes $x = \e^{\I (\pi - \theta)/2} / (2 \sin \tfrac{\theta}{2})$, which lies in the upper half
plane with $\tfrac12 \le \abs{x} \le 1$ while $1 - x$ lies in the lower half plane; with all
branches principal, $w^{-\gamma} = \e^{\I \pi \gamma} x^{-\gamma} (1 - x)^{\gamma}$ and
$(w - 1)^{\gamma - 1} = \e^{-\I \pi (\gamma - 1)} (1 - x)^{1 - \gamma}$, the constant phases
combining with the sign of the Jacobian to $+1$.
The two Pfaff transformations of the ${}_2F_1$ then give two representatives of the same
integrand,
\begin{subequations}
  \label{jb:eq:pfaff}
\begin{align}
  \label{jb:eq:pfaff-a}
  I \, \dd w
    &
    = x^{-\gamma} (1 - x)^{\alpha - 1} \,
      {}_2F_1(\alpha, \alpha + \gamma; \alpha + 1; x) \, \dd x,
  \\
  \label{jb:eq:pfaff-b}
  I \, \dd w
    &
    = x^{-\gamma} (1 - x)^{-\gamma} \,
      {}_2F_1(1, 1 - \gamma; \alpha + 1; x) \, \dd x,
\end{align}
\end{subequations}
and these must not be mixed: the powers of one against the hypergeometric function of the other
is neither.
On \eqref{jb:eq:pfaff-a} the condition \eqref{jb:eq:euler-condition} reads
$\alpha - 1 = \alpha + (\alpha + \gamma) - (\alpha + 1)$, that is $\gamma = 0$, that is $q = -2$.
On \eqref{jb:eq:pfaff-b} it reads $-\gamma = 1 + (1 - \gamma) - (\alpha + 1)$, that is $\alpha = 1$,
that is $p = 0$.
Applied directly to $I$ without the Pfaff step, where the power is $(1 - w)^{\gamma - 1}$ up to a
constant phase, it reads $\gamma - 1 = \alpha + (1 - \gamma) - (\alpha + 1)$, that is
$\gamma = \tfrac12$, that is $q = -1$.
So Euler absorption reduces the circle integrand to a single hypergeometric series on the three
lines
\begin{equation}
  \label{jb:eq:circle-lines}
  p
    = 0,
  \qquad
  q
    = -1,
  \qquad
  q
    = -2,
\end{equation}
and there is a fourth line, $p + q = -4$, that is $\gamma = - \alpha$, where the ${}_2F_1$ itself
degenerates, ${}_2F_1(\alpha, 1 + \alpha; \alpha + 1; w) = (1 - w)^{-\alpha}$, so that
$G = \alpha^{-1} (Z / (1 - Z))^{\alpha} + \text{const}$ is elementary and every arc integrand is
algebraic: an incomplete Beta function on the circle and on the vertical segment, and on the
ellipse a product of four linear factors in the uniformizing variable of
\cref{jb:sec:obstruction:closure}, a Lauricella $F_D^{(3)}$ rather than an Appell function.
Each of these lines closes the circle arc only; that $J$ closes on them is not claimed: on
$q = -1$ and $q = -2$ the ellipse arc is untouched, except at the point $(-1, -1)$ below.

The line $q = -1$ is worth making explicit, since it is the one that will close a point.
With $\gamma = \tfrac12$, Euler's transformation on $I$ gives
$I \, \dd w = \I \, w^{-1/2} \, {}_2F_1(1, \alpha + \tfrac12; \alpha + 1; w) \, \dd w$, the phase
$\I$ coming from $(w - 1)^{-1/2} (1 - w)^{1/2}$ on the lower arc, and \eqref{jb:eq:one-power} then
integrates it:
\begin{equation}
  \label{jb:eq:circle-q-minus-one}
  \int I \, \dd w
    = 2 \I \, w^{1/2} \,
      {}_3F_2 \bigl( 1, \alpha + \tfrac12, \tfrac12; \alpha + 1, \tfrac32; w \bigr)
      \Big|_{-1}^{\e^{- \I \pi/3}},
\end{equation}
the value at $-1$ being taken from below, $w^{1/2} = - \I$.
Both endpoints lie on $\abs{w} = 1$, where this ${}_3F_2$ series converges absolutely for every
$\alpha$, its parameter excess being
$(\alpha + 1 + \tfrac32) - (1 + \alpha + \tfrac12 + \tfrac12) = \tfrac12 > 0$; the endpoint values
are values of the series itself, not analytic continuations.
Against quadrature of $I$ along the arc at $p \in \{-3.4, 1.7, 0.4\}$,
\eqref{jb:eq:circle-q-minus-one} agrees to $3.2 \times 10^{-17}$ on values of size $1.4$ to $2.6$.

\subparagraph{The point $(-1, -1)$.}
Put together, $J(-1, -1)$ is a closed form.
The real segment gives $0$; the vertical segment is on $p = q$ and gives \eqref{jb:eq:J-vert}; the
ellipse is at its one constant-phase point and gives the triangle \eqref{jb:eq:J-ellipse-point}; and
the circle is on $q = -1$, so its swept area is the real part of \eqref{jb:eq:circle-reduced} with the
integral evaluated by \eqref{jb:eq:circle-q-minus-one}.
The result is a combination of ${}_2F_1$ endpoint values and one ${}_3F_2$.
Through \eqref{jb:eq:pq-mandelstam} the point is a codimension-two locus in the ordered Mandelstam
triple, so it has no bearing on the application;
it is recorded because it is the only point off the lines $p = 0$, $q = 0$ where the boundary
route closes.

The conclusion is therefore about the \emph{exponents}, not about the geometry of the region and
not about the endpoints.
For generic $(p, q)$ the arc integrals are (two powers) $\times$ ${}_2F_1$ with exponents unmatched
to the hypergeometric parameters, and no antiderivative in the classes considered was found for
them.
The elliptic arc gives the same verdict, with two extra complications.
Its Schwarz function is two-valued.
And in the uniformizing variable $w$ of \cref{jb:sec:obstruction:closure}, with $c = 4 \beta^2$ and
$x = w / c$, the factor $f \, \dd Z$ is a product of powers of linear factors while
$\overline{G} = G(\bar Z) + \text{const}$ places the ${}_2F_1$ of \eqref{jb:eq:incomplete-beta} at the
argument $\bar Z = (1 + x)^2 / 4 x$ -- the quadratic-rational argument of Gauss's quadratic
transformations.
Pulling the hypergeometric equation back through $z = (1 + x)^2 / 4x$, which is ramified over
$z = 0$ and $z = 1$ only, gives a Fuchsian equation in $x$ with four singular points; it is again a
three-point equation, so that the ${}_2F_1$ in $x$ could be reduced by a quadratic transformation,
only when one ramified point becomes ordinary, which for the parameters
$(\alpha, 1 - \gamma; \alpha + 1)$ means $p \in \{-1, -3\}$ or $q \in \{-1, -3\}$.
This is an argument about where a reduction of the ellipse arc could come from, not a checked
result; what it says is that the ellipse arc is more obstructed than the circle, not equally.

\subparagraph{Verdict of the boundary route.}
For generic $(p, q)$, $J(p, q)$ is an exact one-dimensional integral of a Gauss ${}_2F_1$ times its
derivative, and no closed form for it was found by this route.
Two caveats attach to this statement and should travel with it.
It is a statement about \emph{this route}: it says that reducing each arc contribution to endpoint
values of ${}_2F_1$ and Appell functions does not work, not that no closed form for $J$ exists.
And it is a statement about \emph{generic} $(p, q)$: the next subsection exhibits two lines where
the route does close, and the point $(-1, -1)$ above is a third locus.
The negative part of it is not machine-checkable, since what it asserts is that a route fails,
not that no antiderivative exists; what is machine-checked are its two ingredients, the failure of
the constant-phase mechanism off the loci of \cref{jb:sec:arcs} and the validity of
\eqref{jb:eq:reflection}.


\subsubsection{The lines $p = 0$ and $q = 0$, where the route does close}
\label{jb:sec:obstruction:closure}

On each of the two lines one of the two powers in $f$ is absent, $G$ is elementary, and every arc
of $\partial R$ closes: the real segment gives zero as always, the vertical segment gives an
incomplete Beta function, the unit circle gives an incomplete Beta function on $p = 0$ and an
elementary term on $q = 0$, and the elliptic quarter gives two Appell $F_1$ functions once it is
written in the uniformizing variable of the vertex boundary.
We give every coefficient, phase and branch, one line at a time, and then the checks.
The two lines are treated separately because no symmetry of $R$ maps one to the other; the reason
is spelled out at the end of the subsection.

\subparagraph{Notation.}
Write $c := 4 \beta^2$, so that $a = \tfrac{c^2 + 1}{4c}$, $b = \tfrac{c^2 - 1}{4c}$, and
$c \ge 3 + 2\sqrt2 \approx 5.83$ for $\beta \ge \beta_{04}$ ($c = 9$ at the test value $\beta = \tfrac32$).
The elliptic quarter $C_4$ is parametrized by the uniformizer of the vertex boundary,
\begin{equation}
  \label{jb:eq:uniformizer}
  Z
    = \frac{(c w + 1)^2}{4 c w}
    = \frac12 + \beta^2 w + \frac{1}{16 \beta^2 w},
  \qquad
  w
    = \e^{\I t},
  \quad
  t : \tfrac{\pi}{2} \to \pi,
\end{equation}
so that $w$ runs from $\I$ (at $Z = \tfrac12 + \I b$) to $-1$ (at $Z = \tfrac12 - a$), in the
orientation of $\partial R$.
On the arc,
\begin{equation}
  \label{jb:eq:uniformizer-aux}
  \begin{aligned}
    1 - Z
      &
      = - \frac{(c w - 1)^2}{4 c w},
    &
    \bar Z
      &
      = \frac{(w + c)^2}{4 c w},
    \\
    1 - \bar Z
      &
      = - \frac{(c - w)^2}{4 c w},
    &
    \dd Z
      &
      = \frac{(c w - 1)(c w + 1)}{4 c w^2} \, \dd w,
  \end{aligned}
\end{equation}
and powers of $w$ are taken with $\arg w = t \in [\tfrac{\pi}{2}, \pi]$, so $w^s = \e^{\I \pi s}$
at $w = -1$.
The Appell function enters through the one-parameter family for which it is an antiderivative in
its own base point: with third parameter $a + 1$, the Euler integral \cite{Nist:DLMF} reads
\begin{equation}
  \label{jb:eq:appell-antiderivative}
  \mc F_a(w; b_1, b_2; x, y)
    := \frac{w^a}{a} \, F_1(a; b_1, b_2; a + 1; x, y)
    = \int_0^w s^{a - 1}
      \Bigl( 1 - \frac{x}{w} \, s \Bigr)^{-b_1}
      \Bigl( 1 - \frac{y}{w} \, s \Bigr)^{-b_2}
      \dd s ,
\end{equation}
valid for $\Re a > 0$ and continued analytically in $a$ (simple poles at $a \in \Z_{\le 0}$, which cancel in every endpoint
difference below) and in $x$ off the cut $[1, \infty)$.
Below, $x$ is $\pm c w$ and $y$ is $\pm w / c$: the first argument has $\abs{x} = c \ge 3 + 2\sqrt2$, outside
the unit disk, and needs the continuation of $F_1$; the second has $\abs{y} \le 3 - 2\sqrt2 \approx 0.17$
($\tfrac19$ at $\beta = \tfrac32$).
A form convenient for evaluation follows by expanding the last factor of
\eqref{jb:eq:appell-antiderivative} in $y$,
\begin{equation}
  \label{jb:eq:appell-single-sum}
  F_1(a; b_1, b_2; a + 1; x, y)
    = \sum_{m \ge 0} \frac{(a)_m (b_2)_m}{(a + 1)_m \, m!} \, y^m \,
      {}_2F_1(a + m, b_1; a + m + 1; x),
\end{equation}
a geometrically convergent series in $y$ whose terms need only a ${}_2F_1$ at $\abs{x} = c$.

\subparagraph{The line $p = 0$.}
Here $f = (1 - Z)^{q/2}$ and, for $\gamma \ne 0$, the antiderivative normalized by
\eqref{jb:eq:base-point} is
\begin{equation}
  \label{jb:eq:G-p-zero}
  G(Z)
    = \frac{2^{\gamma} - (1 - Z)^{\gamma}}{\gamma},
  \qquad
  \gamma
    = \frac{q}{2} + 1 .
\end{equation}
Since $\overline{(1 - Z)^{\gamma}} \, (1 - Z)^{q/2} = \abs{1 - Z}^{q + 1} \e^{-\I \theta_1}
= \abs{1 - Z}^{q} (1 - \bar Z)$, the swept-area integrand splits into an exact part and a part
without $G$,
\begin{equation}
  \label{jb:eq:reduction-p-zero}
  \overline{G} \, \dd G
    = \frac{2^{\gamma}}{\gamma} \, \dd G
      - \frac{1}{\gamma} \, \abs{1 - Z}^{q} \, (1 - \bar Z) \, \dd Z ,
\end{equation}
so that per arc
\begin{equation}
  \label{jb:eq:per-arc-p-zero}
  \begin{aligned}
    J_k
      &
      = \frac{2^{\gamma}}{2 \gamma} \, \Im \bigl[ G \bigr]_k
        - \frac{1}{2 \gamma} \, \Im I_k,
    \qquad
    I_k
      := \int_k \abs{1 - Z}^{q} \, (1 - \bar Z) \, \dd Z,
    \\
    J(0, q)
      &
      = - \frac{1}{q + 2} \, \Im \sum_k I_k,
  \end{aligned}
\end{equation}
the endpoint terms $[G]_k$ telescoping around the closed boundary.
(They are kept in the per-arc formula so that each $J_k$ can be compared separately with the swept
area of \cref{jb:sec:stokes} evaluated by quadrature; the totals do not need them.)
What remains is to write $\abs{1 - Z}^{q} (1 - \bar Z)$ holomorphically on each arc through its
Schwarz function \eqref{jb:eq:schwarz}, and at each step the point to check is that the principal
arguments of the factors add up to the argument of the product without a $2\pi$ defect.

On the unit circle, $\bar Z = 1/Z$ gives $\abs{1 - Z}^2 = (1 - Z)(1 - 1/Z)$.
With $\theta_1 = \tfrac{\theta - \pi}{2}$ one has $\arg(Z - 1) - \arg Z = \tfrac{\pi - \theta}{2}
\in [0, \tfrac{\pi}{3}]$, so the principal powers compose, $(1 - 1/Z)^{q/2} = \e^{\I \pi q/2}
Z^{-q/2} (1 - Z)^{q/2}$, and
\begin{equation}
  \label{jb:eq:closure-circle}
  \begin{aligned}
    \abs{1 - Z}^{q} \, (1 - \bar Z) \, \dd Z
      &
      = - \e^{\I \pi q/2} \, Z^{-q/2 - 1} (1 - Z)^{q + 1} \, \dd Z,
    \\
    I_{\text{circle}}
      &
      = - \e^{\I \pi q/2} \Bigl[ B_Z \bigl( - \tfrac{q}{2}, \ q + 2 \bigr) \Bigr]_{-1}^{\e^{\I \pi/3}} .
  \end{aligned}
\end{equation}
On the vertical segment, $\bar Z = 1 - Z$ and $\theta_1 = - \theta_0$ give
$\abs{1 - Z}^{q} = Z^{q/2} (1 - Z)^{q/2}$ exactly, hence
\begin{equation}
  \label{jb:eq:closure-vert}
  \begin{aligned}
    \abs{1 - Z}^{q} \, (1 - \bar Z) \, \dd Z
      &
      = Z^{q/2 + 1} (1 - Z)^{q/2} \, \dd Z,
    \\
    I_{\text{vert}}
      &
      = \Bigl[ B_Z \bigl( \tfrac{q}{2} + 2, \ \tfrac{q}{2} + 1 \bigr) \Bigr]_{\e^{\I \pi/3}}^{\frac12 + \I b} .
  \end{aligned}
\end{equation}
On the real segment the integrand $(1 - x)^{q + 1} \dd x$ and $G(x)$ are real, and
$J_{\text{real}} = 0$ as in \cref{jb:sec:arcs:real}.

The elliptic quarter is the only arc needing work.
Put $A := c w - 1$ and $B := c - w$; on $\abs{w} = 1$ one has $\bar A = B / w$, so the positive
real number $\abs{1 - Z} = \abs{A}^2 / 4c$ is the rational function $A B / 4 c w$ of $w$, and the
principal arguments satisfy $\arg A + \arg B - \arg w = 0$ exactly on the arc ($\arg A \in
(\tfrac{\pi}{2}, \pi]$, $\arg B = t - \arg A$, checked at both ends).
Therefore
\begin{equation}
  \label{jb:eq:closure-ellipse-integrand}
  \abs{1 - Z}^{q} \, (1 - \bar Z) \, \dd Z
    = - (4 c)^{-q - 2} \, A^{q + 1} \, B^{q + 2} \, (c w + 1) \, w^{-q - 3} \, \dd w .
\end{equation}
To bring this to the Appell form \eqref{jb:eq:appell-antiderivative}, note that $1 - c w = - A$ has
argument in $(- \tfrac{\pi}{2}, 0]$, so $A^{q + 1} = \e^{\I \pi (q + 1)} (1 - c w)^{q + 1}$ with
both powers principal, and $B^{q + 2} = c^{q + 2} (1 - w / c)^{q + 2}$ with $\abs{w / c} \le
\tfrac19$.
The constants combine into $- (4 c)^{-q - 2} c^{q + 2} \e^{\I \pi (q + 1)} = 4^{-q - 2} \e^{\I \pi q}$,
and the linear factor splits as $(c w + 1) w^{-q - 3} = w^{-q - 3} + c \, w^{-q - 2}$, two terms of
the shape $w^{a - 1}$ with $a = -q - 2$ and $a = -q - 1$:
\begin{equation}
  \label{jb:eq:closure-ellipse}
  \begin{aligned}
    I_{\text{ell}}
      &
      = 4^{-q - 2} \, \e^{\I \pi q}
        \Bigl[ \mc F_{-q - 2}(w) + c \, \mc F_{-q - 1}(w) \Bigr]_{w = \I}^{w = -1},
    \\
    \mc F_a(w)
      &
      := \frac{w^a}{a} \, F_1 \bigl( a; \ - q - 1, \ - q - 2; \ a + 1; \ c w, \ \tfrac{w}{c} \bigr) .
  \end{aligned}
\end{equation}
That the endpoint difference of the straight-path integrals \eqref{jb:eq:appell-antiderivative}
equals the integral along the arc is Cauchy's theorem on the quarter disk $\{ \Re w \le 0, \
\Im w \ge 0, \ \abs{w} \le 1 \}$ for $\Re a > 0$: the integrand is holomorphic inside, the cuts of
$(1 - c w)^{q + 1}$ and $(1 - w / c)^{q + 2}$ lie on $[1/c, \infty)$ and $[c, \infty)$, outside the
disk, and $w^{a - 1}$ carries its boundary value from above on the negative axis, where the path
$0 \to -1$ runs.
For all other $a$ the identity extends by analytic continuation, both sides being meromorphic in
$a$ with the same poles.
The first Appell argument $x = c w \in \{ c \I, - c \}$ stays off the cut of $F_1$.

Summing the three arcs, with the endpoint terms of \eqref{jb:eq:per-arc-p-zero} telescoping,
\begin{equation}
  \label{jb:eq:J-p-zero}
  \boxed{\
  \begin{aligned}
    J(0, q)
      &
      = - \frac{1}{q + 2} \, \Im \Bigl\{
        - \e^{\I \pi q/2} \Bigl[ B_Z \bigl( - \tfrac{q}{2}, \ q + 2 \bigr) \Bigr]_{-1}^{\e^{\I \pi/3}}
        + \Bigl[ B_Z \bigl( \tfrac{q}{2} + 2, \ \tfrac{q}{2} + 1 \bigr) \Bigr]_{\e^{\I \pi/3}}^{\frac12 + \I b}
      \\
      &
      \qquad\qquad\qquad
        + 4^{-q - 2} \, \e^{\I \pi q}
          \Bigl[ \mc F_{-q - 2} + c \, \mc F_{-q - 1} \Bigr]_{\I}^{-1}
      \Bigr\}
  \end{aligned}
  \ }
\end{equation}
valid as written for $q \notin \{ 0, -2 \} \cup \{ -4, -6, -8, \dots \} \cup \Z_{\ge -1}$: the
normalization $Z^{\alpha'} / \alpha'$ of the circle term has a pole at $q \in \{ 0, 2, 4, \dots \}$
and that of the vertical term at $q \in \{ -4, -6, \dots \}$, the prefactor has one at $q = -2$,
and $F_1$ has third parameter $a + 1 \in \Z_{\le 0}$ at integer $q \ge -1$.
All of these are removable, $J$ being entire in $q$ (\cref{jb:sec:setup}), and the formula extends to
the excluded values by continuity.

\subparagraph{The line $q = 0$.}
Here $f = Z^{p/2}$ and, for $\alpha \ne 0$,
\begin{equation}
  \label{jb:eq:G-q-zero}
  G(Z)
    = \frac{Z^{\alpha} - \e^{\I \pi \alpha}}{\alpha},
  \qquad
  \alpha
    = \frac{p}{2} + 1,
\end{equation}
with $(-1)^{\alpha} = \e^{\I \pi \alpha}$ the boundary value from above of \eqref{jb:eq:branches}.
Now $\overline{Z^{\alpha}} \, Z^{p/2} = \abs{Z}^{p + 1} \e^{-\I \theta_0} = \abs{Z}^{p} \bar Z$, and
\begin{equation}
  \label{jb:eq:reduction-q-zero}
  \begin{aligned}
    \overline{G} \, \dd G
      &
      = \frac{1}{\alpha} \, \abs{Z}^{p} \, \bar Z \, \dd Z
        - \frac{\e^{-\I \pi \alpha}}{\alpha} \, \dd G,
    \\
    J_k
      &
      = \frac{1}{2 \alpha} \, \Im K_k
        - \frac{1}{2 \alpha^2} \, \Im \Bigl( \e^{-\I \pi \alpha} \bigl[ Z^{\alpha} \bigr]_k \Bigr),
    \qquad
    K_k
      := \int_k \abs{Z}^{p} \, \bar Z \, \dd Z,
  \end{aligned}
\end{equation}
so that $J(p, 0) = \tfrac{1}{2 \alpha} \Im \sum_k K_k$, the endpoint terms telescoping because
$Z^{\alpha}$ is single-valued on $\partial R$ with the boundary value from above.
(The endpoint term carries $\alpha^{-2}$, from $\dd G = \dd Z^{\alpha} / \alpha$; writing
$\alpha^{-1}$ there leaves the total intact but breaks the per-arc comparison.)

On the unit circle $\abs{Z} = 1$ and $\bar Z \, \dd Z = Z^{-1} \dd Z$, so the arc is elementary,
with no hypergeometric function at all:
\begin{equation}
  \label{jb:eq:closure-circle-q-zero}
  K_{\text{circle}}
    = \bigl[ \ln Z \bigr]_{-1}^{\e^{\I \pi/3}}
    = - \frac{2 \pi \I}{3} .
\end{equation}
(The incomplete Beta function one would write by analogy with \eqref{jb:eq:closure-circle}
degenerates: its exponents are $(-1, 0)$.)
On the vertical segment, $\abs{Z}^{p} = Z^{p/2} (1 - Z)^{p/2}$ exactly and $\bar Z = 1 - Z$, so
\begin{equation}
  \label{jb:eq:closure-vert-q-zero}
  \abs{Z}^{p} \, \bar Z \, \dd Z
    = Z^{p/2} (1 - Z)^{p/2 + 1} \, \dd Z,
  \qquad
  K_{\text{vert}}
    = \Bigl[ B_Z (\alpha, \ \alpha + 1) \Bigr]_{\e^{\I \pi/3}}^{\frac12 + \I b} .
\end{equation}
On the real segment $\abs{x}^{p} x \, \dd x$ and $\e^{-\I \pi \alpha} Z^{\alpha} = \abs{x}^{\alpha}$
are real, and $J_{\text{real}} = 0$.

On the elliptic quarter, $\overline{c w + 1} = (c + w) / w$ gives $\abs{Z} = (c w + 1)(w + c) /
4 c w$, the principal arguments composing exactly ($\arg(c w + 1) + \arg(w + c) - \arg w = 0$ on
the arc: $\pi + 0 - \pi$ at $w = -1$ and $\arctan c + \arctan \tfrac1c - \tfrac{\pi}{2}$ at
$w = \I$).
With $(w + c)^{p + 2} = c^{p + 2} (1 + w / c)^{p + 2}$, $c w - 1 = - (1 - c w)$ and
$(4 c)^{-p - 2} c^{p + 2} = 4^{-p - 2}$,
\begin{equation}
  \label{jb:eq:closure-ellipse-integrand-q-zero}
  \begin{aligned}
    \abs{Z}^{p} \, \bar Z \, \dd Z
      &
      = - 4^{-p - 2} \, (1 + c w)^{p + 1} (1 + w / c)^{p + 2} (1 - c w) \, w^{-p - 3} \, \dd w
      \\
      &
      = - 4^{-p - 2} \bigl[ w^{-p - 3} - c \, w^{-p - 2} \bigr]
        (1 + c w)^{p + 1} (1 + w / c)^{p + 2} \, \dd w,
  \end{aligned}
\end{equation}
and therefore
\begin{equation}
  \label{jb:eq:closure-ellipse-q-zero}
  \begin{aligned}
    K_{\text{ell}}
      &
      = - 4^{-p - 2}
        \Bigl[ \widetilde{\mc F}_{-p - 2}(w) - c \, \widetilde{\mc F}_{-p - 1}(w) \Bigr]_{w = \I}^{w = -1},
    \\
    \widetilde{\mc F}_a(w)
      &
      := \frac{w^a}{a} \, F_1 \bigl( a; \ - p - 1, \ - p - 2; \ a + 1; \ - c w, \ - \tfrac{w}{c} \bigr) .
  \end{aligned}
\end{equation}
One prescription is needed here that the line $p = 0$ did not need.
The power $(1 + c w)^{p + 1}$ is principal with $\Im(1 + c w) = c \sin t \ge 0$, and at $w = -1$ it
is the boundary value from above of the cut of $(1 + c w)^{p + 1}$ at $w = - 1/c$.
This is the cut of $Z^{p/2}$ at $Z = 0$, which $\partial R$ touches along the real segment; the
line $p = 0$ has no such cut, which is why \eqref{jb:eq:closure-ellipse} needed nothing of the kind.
Correspondingly the first Appell argument $x = - c w$ is $c - \I 0$ at $w = -1$ -- on the cut
$[1, \infty)$ of $F_1$, approached from $\Im x < 0$ -- and $- c \I$ at $w = \I$.
The Cauchy argument given for \eqref{jb:eq:closure-ellipse} applies with the branch point $w = - 1/c$
on the boundary of the quarter disk; the identity holds for $\Re a > 0$, $\Re (p + 1) > -1$, and
extends analytically in $(a, b_1)$ to all $p$.
A prescription-free form, used by the second evaluator below, follows from substituting
$u = 1 + c w$ in each term of \eqref{jb:eq:appell-single-sum}:
\begin{equation}
  \label{jb:eq:closure-ellipse-q-zero-u}
  \int_{\I}^{-1} w^{a + m - 1} (1 + c w)^{p + 1} \, \dd w
    = c^{-a - m} \, \e^{\I \pi (a + m - 1)}
      \Bigl[ B_u (p + 2, \ a + m) \Bigr]_{u = 1 + c \I}^{u = 1 - c},
\end{equation}
where $u^{p + 2}$ at $u = 1 - c < 0$ carries the boundary value from above and the ${}_2F_1$
argument $1 - c$ is off its cut.

Summed,
\begin{equation}
  \label{jb:eq:J-q-zero}
  \boxed{\
  \begin{aligned}
    J(p, 0)
      &
      = \frac{1}{p + 2} \Bigl\{
        - \frac{2 \pi}{3}
        + \Im \Bigl[ B_Z (\alpha, \ \alpha + 1) \Bigr]_{\e^{\I \pi/3}}^{\frac12 + \I b}
      \\
      &
      \qquad\qquad
        - 4^{-p - 2} \, \Im \Bigl[ \widetilde{\mc F}_{-p - 2} - c \, \widetilde{\mc F}_{-p - 1} \Bigr]_{\I}^{-1}
      \Bigr\},
    \qquad
    \alpha
      = \frac{p}{2} + 1,
  \end{aligned}
  \ }
\end{equation}
valid as written for $p \notin \{ -2, -4, -6, \dots \} \cup \Z_{\ge -1}$ (the prefactor at $p = -2$
and the normalization $Z^{\alpha} / \alpha$ of the vertical term at $p \in \{ -2, -4, \dots \}$, as for
\eqref{jb:eq:J-p-zero}; $F_1$ with $a + 1 \in \Z_{\le 0}$ at integer $p \ge -1$), the poles removable and
the formula extended by continuity.

\subparagraph{The limits $q \to 0$ and $p \to 0$.}
At $q = 0$ the three arc integrals of \eqref{jb:eq:J-p-zero} are elementary and
\eqref{jb:eq:per-arc-p-zero} reads $J(0, 0) = \tfrac12 \Im \oint \bar Z \, \dd Z$, the area of $R$:
\eqref{jb:eq:J-p-zero} reproduces the area of $R$ in the limit, and so does \eqref{jb:eq:J-q-zero} at
$p = 0$, where \eqref{jb:eq:closure-circle-q-zero} is the circle's share of the area.
Numerically the poles of $B_Z$ and $F_1$ at these points cancel between endpoints rather than
term by term, so the limit is taken symmetrically, $\tfrac12 [J(0, \epsilon) + J(0, - \epsilon)]
= J(0, 0) + O(\epsilon^2)$.

\subparagraph{The two lines are not mirror images.}
The map $Z \mapsto 1 - \bar Z$ exchanges $\abs{Z}$ and $\abs{1 - Z}$ and preserves the full
confocal ellipse, but it sends $R$ -- which has $\abs{Z} \ge 1$ and $\Re Z \le \tfrac12$ -- to
$\{ \abs{1 - Z} \ge 1, \ \Re Z \ge \tfrac12 \} \cap \mc V_{0,4} \ne R$.
The symmetry group of $R$ is trivial, no element of it relates the two lines, and indeed
$J(0, e) \ne J(e, 0)$: at $\beta = \tfrac32$, $J(0, 1.3) = 7.2777$ while $J(1.3, 0) = 4.7719$.
The structural asymmetry is the cut: $\partial R$ touches the cut of $Z^{p/2}$ along the real
segment and never that of $(1 - Z)^{q/2}$.
This is where the $c - \I 0$ prescription of \eqref{jb:eq:closure-ellipse-q-zero} comes from, and its
absence in \eqref{jb:eq:closure-ellipse}, and it is why the circle is elementary on $q = 0$
(\eqref{jb:eq:closure-circle-q-zero}, $\abs{Z}^p \equiv 1$ there) and an incomplete Beta function on
$p = 0$.
Note that on $q = 0$ the circle does not close by the constant-phase mechanism of
\cref{jb:sec:arcs:curved} -- the line $2p + q = -4$ meets $q = 0$ only at the excluded point
$(-2, 0)$ -- but for this other reason, which is also why $J(p, 0)$ is only algebraically
suppressed at large $\abs{p}$.

\subparagraph{Checks.}
Both closed forms were compared with the boundary representation \eqref{jb:eq:stokes} evaluated arc
by arc by quadrature, at sixteen points: $q$ or $p$ in $\{ -21.7, -3.3, 1.3, 20.3 \}$ on each
line, $\beta \in \{ \tfrac32, 2 \}$.
The worst relative difference is $1.4 \times 10^{-15}$ on the totals and $1.1 \times 10^{-15}$ arc
by arc, which is the double-precision floor of the reference's arc parametrizations.
Multiplying the largest arc term by $1 + 10^{-3}$ makes the comparison fail at all sixteen points,
and so does multiplying the ellipse term at the twelve points where it carries more than $10^{-4}$
of $J$ (at $q = -21.7$ it carries about $10^{-9}$ of it, the weight sitting at $\e^{\I \pi/3}$), so
the check is sensitive to each piece.
The Appell functions were evaluated in two ways sharing no code beyond the ${}_2F_1$: a library
continuation of $F_1$ in its first argument, with $c - \I 0$ realized as $c (1 - 10^{-40} \I)$,
and the single-sum form \eqref{jb:eq:appell-single-sum} -- on $q = 0$ in its incomplete-Beta version
\eqref{jb:eq:closure-ellipse-q-zero-u}.
They agree to $1.2 \times 10^{-26}$ at all sixteen points, which fixes both the continuation and
the side of the cut: taking $c + \I 0$ instead changes the result at order one (relative
difference $0.94$).
The symmetric limits above reproduce the exact area of $R$ to $6.8 \times 10^{-13}$ at
$\epsilon = 10^{-6}$ (a one-sided limit is off by $1.1 \times 10^{-6}$, as it should be).
Two things are not checked: the formulas at the excluded integer exponents, which are taken by
continuity, and the continuation of $F_1$ beyond the mutual agreement of the two evaluators -- no
third, direct evaluation of the Euler integral at $\Re a \le 0$ was run.

\subparagraph{What the closure means.}
It does not help the application: through \eqref{jb:eq:pq-mandelstam}, $p = 0$ and $q = 0$ are the
codimension-one loci $x = y + z$ and $y = x + z$ in the ordered Mandelstam triple, a measure-zero
subset of any grid one would want to scan.
But the two lines lie inside the class the negative statement of \cref{jb:sec:obstruction:powers}
covers, which is why that statement must be qualified as holding for generic $(p, q)$, and they
give two exact anchors against which a series representation could be calibrated.



\subsubsection{A lead not followed}
\label{jb:sec:obstruction:lead}

One consequence of \eqref{jb:eq:reflection} deserves recording, since it may be worth more than it
looks.
Because the circle integrand is holomorphic after the reflection, the contour may be deformed off
$\abs{Z} = 1$ into $\abs{Z} > 1$, staying clear of $Z = 1$.
This is exactly what an expansion in $1/Z$ would want, since the expansion is marginal on the unit
circle itself.
The catch is that the endpoints $-1$ and $\e^{\I \pi/3}$ are fixed and stay on $\abs{Z} = 1$, with
$-1$ additionally sitting on the cut of $B_{1/Z}$.
The deformation therefore improves the interior of the contour and not its ends, and whether it
rescues a $1/Z$ expansion of this piece is not examined here.


\subsection{The confocal structure of the vertex region}
\label{jb:sec:confocal}


This section is independent of everything above.
It records a property of the vertex region itself which, once seen, changes which approximation
schemes look promising -- and which identifies the boundary curve responsible for the difficulty.

\subparagraph{The vertex boundary is a confocal ellipse.}
Start from $a^2 - b^2 = \tfrac14$ in \eqref{jb:eq:axes}, equivalently $c = \tfrac12$ for the focal
distance of the $s$-channel ellipse, which is centered at $z_c = \tfrac12$.
Its foci therefore sit at $\tfrac12 \pm \tfrac12$, that is at $0$ and $1$: exactly the two
punctures, exactly the two branch points of the integrand.
By the focal definition of an ellipse,
\begin{equation}
  \label{jb:eq:confocal}
  \pd \mc V_{0,4}
    = \Acom{ \abs{Z} + \abs{1 - Z} = 2a } .
\end{equation}
The vertex boundary is a level set of the sum of the distances to the punctures.
This is not an approximation valid at large stub length; it is exact for every $\beta$.

\subparagraph{Bipolar coordinates.}
The natural coordinates are then the two distances themselves,
\begin{equation}
  \label{jb:eq:uv-def}
  r
    := \abs{Z},
  \qquad
  \rho
    := \abs{1 - Z},
  \qquad
  u
    := r + \rho,
  \qquad
  v
    := \rho - r .
\end{equation}
From $\rho^2 = 1 - 2x + r^2$ one gets $x = \tfrac12 (r^2 - \rho^2 + 1)$ and $y = \sqrt{r^2 - x^2}$ on
the sheet $y \ge 0$, so that
\begin{equation}
  \label{jb:eq:jacobian-1}
  \frac{\pd (x, y)}{\pd (r, \rho)}
    = \frac{r \rho}{y},
  \qquad
  4 y^2
    = \bigl( (r + \rho)^2 - 1 \bigr) \bigl( 1 - (\rho - r)^2 \bigr)
    = (u^2 - 1)(1 - v^2).
\end{equation}
The factorization of $4y^2$ is not an accident: $y$ is twice the area of the triangle with vertices
$0$, $1$, $Z$, whose sides are $1$, $r$, $\rho$, and the middle expression in
\eqref{jb:eq:jacobian-1} is Heron's formula for it.
Using $\dd r \, \dd \rho = \tfrac12 \dd u \, \dd v$ and $r \rho = \tfrac14 (u^2 - v^2)$,
\begin{equation}
  \label{jb:eq:jacobian-2}
  \dd x \, \dd y
    = \frac{(u^2 - v^2) \, \dd u \, \dd v}{4 \sqrt{(u^2 - 1)(1 - v^2)}}
  \qquad (y \ge 0).
\end{equation}

\subparagraph{The region becomes a trapezoid.}
Each of the four conditions \eqref{jb:eq:region} is now simple.
The vertex boundary is $u = 2a$ by \eqref{jb:eq:confocal}; the condition $\Re Z \le \tfrac12$ says that
$Z$ is at least as far from $1$ as from $0$, that is $v \ge 0$; the condition $\abs{Z} \ge 1$ reads
$u - v \ge 2$; and the triangle inequality on $0$, $1$, $Z$ gives $v \le 1$ automatically, with
equality on the negative real axis.
Hence
\begin{equation}
  \label{jb:eq:trapezoid}
  R
    \ \longleftrightarrow \
    D
    = \Acom{ 0 \le v \le 1, \ 2 + v \le u \le 2a },
\end{equation}
and \eqref{jb:eq:J-def} becomes
\begin{equation}
  \label{jb:eq:J-uv}
  J(p, q)
    = 2^{-p-q-2} \int_0^1 \frac{\dd v}{\sqrt{1 - v^2}}
      \int_{2+v}^{2a} \frac{\dd u}{\sqrt{u^2 - 1}} \,
      (u - v)^{p+1} (u + v)^{q+1} .
\end{equation}
The four corners of $D$ are $(2, 0) \leftrightarrow \e^{\I \pi/3}$,
$(3, 1) \leftrightarrow -1$, $(2a, 1) \leftrightarrow \tfrac12 - a$ and
$(2a, 0) \leftrightarrow \tfrac12 + \I b$.
The sides $v = 0$, $v = 1$ and $u = 2a$ are respectively the vertical segment, the real segment and
the elliptic arc -- all three are coordinate lines.
The slanted side $u = 2 + v$ is the unit circle, and it is the only one that is not.

\subparagraph{Physical reading.}
The vertex boundary, which is the complicated-looking curve in the $Z$ plane, is the \emph{easy}
side of the region in these coordinates; the unit circle, which is the simple-looking curve, is the
hard one.
The obstruction to separating variables in \eqref{jb:eq:J-uv} is the constraint $\abs{Z} \ge 1$
inherited from the $S_4$ fundamental domain, not the stub-dependent vertex boundary.
Any scheme that tries to compute $J$ by expanding around the region's geometry should be built
around that fact.

\subparagraph{Check: the full Virasoro--Shapiro integral.}
The same Jacobian on the whole plane -- both sheets $y \gtrless 0$, hence $v \in [-1, 1]$ and
$u \ge 1$ -- turns the unrestricted integral into
$2 \cdot 2^{-p-q-2} \int_{-1}^{1} \! \int_1^{\infty}$ of the same integrand, which must reproduce the
closed-form Virasoro--Shapiro result
\begin{equation}
  \label{jb:eq:VS}
  \int_{\C} \abs{Z}^p \abs{1 - Z}^q \, \dd^2 Z
    = \pi \,
      \frac{\Gamma(1 + \tfrac{p}{2}) \, \Gamma(1 + \tfrac{q}{2}) \,
        \Gamma(- 1 - \tfrac{p+q}{2})}
        {\Gamma(- \tfrac{p}{2}) \, \Gamma(- \tfrac{q}{2}) \,
        \Gamma(2 + \tfrac{p+q}{2})}
\end{equation}
in the convergence window $p, q > -2$, $p + q < -2$.
It does, to $4 \times 10^{-11}$ at three test points.
This is the cleanest available test of \eqref{jb:eq:jacobian-2}, since it compares against a formula
derived by completely different means.

\subparagraph{A series, and why it is not finished here.}
In \eqref{jb:eq:J-uv} the ratio $w = v/u = (\rho - r)/(\rho + r)$ satisfies $0 \le w \le \tfrac13$ on
all of $R$, with the maximum attained at $Z = -1$.
Expanding,
\begin{equation}
  \label{jb:eq:series}
  (u - v)^{p+1} (u + v)^{q+1}
    = u^{p+q+2} \sum_{n \ge 0} c_n(p, q) \, w^n,
  \qquad
  c_n
    = \sum_{j=0}^{n} (-1)^j \binom{p+1}{j} \binom{q+1}{n-j},
\end{equation}
which converges geometrically at ratio $\le \tfrac13$, uniformly in $Z$ on $R$.
This is a genuine improvement over an expansion in $1/Z$, which is marginal on the unit circle.
It is, however, \emph{not} uniform in $(p, q)$: the coefficients $c_n$ grow like
$\binom{\abs{p} + n}{n}$, so the geometric decay only sets in after $n \gtrsim \abs{p}/2$, and the
number of terms needed for a given accuracy grows with $\abs{p} + \abs{q}$.

Each term splits over the region in the obvious way.
On the rectangle $[3, 2a] \times [0, 1]$ the two variables separate and the term factorizes into
$\tfrac12 B(\tfrac{n+1}{2}, \tfrac12)$ times $\int_3^{2a} u^{p+q+2-n} (u^2-1)^{-1/2} \dd u$, a
${}_2F_1$ in $1/u^2$.
On the triangle $\{ 2 \le u \le 3, \ 0 \le v \le u - 2 \}$ the lower limit depends on $v$ and the
variables stay coupled.
That triangle is the neighborhood of the unit-circle arc, and it is where a closed form for the
individual terms would have to come from.
Neither the term integrals over the triangle nor a rigorous tail bound for \eqref{jb:eq:series} are
established here.

The representation \eqref{jb:eq:J-uv} has been checked against the boundary representation
\eqref{jb:eq:stokes} and against two-dimensional quadrature at random $(p, q)$, agreeing to
$2 \times 10^{-15}$ with the former and to $5 \times 10^{-11}$ with the latter.


\subsection{Numerical evaluation}
\label{jb:sec:numerics}


The identity \eqref{jb:eq:stokes} is exact, so its numerical value is limited only by how well the four
arc integrals are done.
Two evaluations were implemented, deliberately sharing as little as possible.

\subparagraph{The hypergeometric route.}
Take $G$ from \eqref{jb:eq:G-def}, with the ${}_2F_1$ evaluated at $30$-digit precision, and do each
arc integral by adaptive one-dimensional quadrature.
This is the definitional route and it is accurate to essentially arbitrary precision, but it is
slow: up to $20$ seconds per value of $J$ at $\abs{p} + \abs{q} \sim 170$, dominated by the
hypergeometric evaluations.

\subparagraph{The Chebyshev route, which is the one to use.}
No hypergeometric function is needed at all, because $G$ is only ever required \emph{along the
contour}.
On each arc set $h(t) = f(Z(t)) \, Z'(t)$, which involves only powers and elementary functions.
Interpolate $h$ by a Chebyshev series of degree $600$ in the arc parameter, integrate the
interpolant in closed form -- Chebyshev series have exact term-by-term antiderivatives -- and
accumulate the result from arc to arc starting at $G(-1) = 0$.
This produces $G$ along the whole boundary.
Then take $\tfrac12 \Im \int \overline{G} \, h \, \dd t$ on each arc by Gauss--Legendre quadrature
with the same number of points.
The total cost is about $2400$ elementary evaluations of $f$ per value of $J$, in double precision.

One caveat on the word ``exactly'' above: the cumulative integration of the interpolant is exact,
but the final quadrature is applied to the \emph{product} $\overline{G} h$, whose Chebyshev degree
is roughly twice that of $h$, while Gauss--Legendre with $N$ points is exact only to degree
$2N - 1$.
The tail coefficients are below the working precision, which is why this is harmless, and the
agreement with the hypergeometric route is the evidence for that.

\subparagraph{Accuracy.}
The two routes were compared against each other and against a direct two-dimensional quadrature of
\eqref{jb:eq:J-def}.

\begin{table}[htp]
  \centering
  \begin{tabular}{lll}
    \toprule
    What is compared & Conditions & Agreement \\
    \midrule
    \eqref{jb:eq:stokes} at $p = q = 0$ vs.\ the area of $R$ & both $\beta$, ${}_2F_1$ route & $3 \times 10^{-18}$ \\
    the same, Chebyshev route & both $\beta$ & $10^{-15}$ \\
    \eqref{jb:eq:stokes} (${}_2F_1$) vs.\ 2D quadrature & 12 random $(p,q)$, both $\beta$ & $5 \times 10^{-10}$ \\
    \eqref{jb:eq:stokes} (Chebyshev) vs.\ 2D quadrature & same 12 points & $5 \times 10^{-10}$ \\
    the two evaluations against each other & same 12 points & $6 \times 10^{-11}$ \\
    the same, at moderate exponents & & $10^{-15}$ \\
    \eqref{jb:eq:J-uv} vs.\ \eqref{jb:eq:stokes} & 4 random points & $2 \times 10^{-15}$ \\
    \eqref{jb:eq:VS} vs.\ the bipolar Jacobian on $\C$ & 3 points in the window & $4 \times 10^{-11}$ \\
    \bottomrule
  \end{tabular}
  \caption{
    Worst relative agreement over the test points.
    The twelve random points comprise six per value of $\beta \in \{ \tfrac32, 2 \}$, of which two
    per $\beta$ have $\abs{p} + \abs{q} > 100$, reaching $(p, q) = (-91, -82)$ and $(-86, +82)$;
    the values of $J$ there span $10^{-5}$ to $10^{22}$.
  }
  \label{jb:tab:accuracy}
\end{table}

Two remarks on \cref{jb:tab:accuracy}.

The $5 \times 10^{-10}$ figure in the third and fourth rows is \emph{the reference's} accuracy, not
the representation's.
The two-dimensional quadrature converges only as $n^{-3}$ in its panel count: its inner limits
$\mp \sqrt{1 - y^2}$ put a square-root branch point at the panel endpoint $y = 1$, where both inner
segments lose length like $\sqrt{2(1-y)}$; the outer parametrization $y = b \sin t$ has nonzero
slope there, so the branch point survives in the integration variable.
Gauss--Legendre against an endpoint singularity $x^{s}$ converges as $n^{-2(s+1)}$, which is
$n^{-3}$ at $s = \tfrac12$.
The measured error ratios between $n = 160$, $320$ and $640$ are close to $8$, confirming the
exponent.
At a typical working value $n = 40$ the two-dimensional quadrature carries about $2 \times 10^{-6}$
relative error for moderate negative exponents, and about $10^{-8}$ at $n = 160$.
The boundary representation is the more accurate of the two sides, and the fifth row
($6 \times 10^{-11}$, and $10^{-15}$ at moderate exponents) is the better estimate of what it
delivers.

The comparisons are falsifiable: multiplying the boundary value by $1 + 10^{-3}$ makes the
agreement test fail, as it must, and the reflection identity \eqref{jb:eq:reflection} fails when an arc
is fed the Schwarz function of a different arc.

\subparagraph{What is not established.}
The error of the Chebyshev route decays geometrically, because $f(Z(t)) Z'(t)$ is analytic in a
neighborhood of each closed arc parameter interval -- $R$ avoids both punctures -- so the standard
Bernstein-ellipse estimate applies \cite{Trefethen:2013:ApproximationTheoryPractice}.
The Bernstein parameter of each arc is computable, but it has not been computed, so there is
\emph{no rigorous error bound} here; what is available is the measured figure of
\cref{jb:tab:accuracy}.
Separately, the bipolar check \eqref{jb:eq:VS} excludes the point $(p, q) = (-0.5, -1.9)$, where the
corner singularity $r^{-0.8}$ at $Z = 1$ leaves the nested tanh--sinh quadrature at
$4 \times 10^{-3}$; that is a limitation of the check, not of \eqref{jb:eq:J-uv}.


\subsection{How to use this}
\label{jb:sec:use}


For any real $(p, q)$ and any $\beta \ge \beta_{04}$, evaluate
\begin{equation*}
  J(p, q)
    = \frac12 \, \Im \oint_{\pd R} \overline{G} \, \dd G
\end{equation*}
by walking once around the four arcs of $\partial R$, accumulating $G$ from $G(-1) = 0$ by
cumulative Chebyshev integration of $f(Z(t)) Z'(t)$, and taking the swept area on each arc.
Degree $600$ suffices; the cost is a few thousand evaluations of elementary functions in double
precision, milliseconds in practice, and the accuracy is about $10^{-10}$ in the worst tested case
and $10^{-15}$ at moderate exponents.
Three points of implementation matter and are easy to get wrong:

\begin{itemize}
  \setlength{\itemsep}{0.5ex}

  \item the branch prescription \eqref{jb:eq:branches}, in particular the boundary value \emph{from
    above} on the negative real axis, which is where the region touches the cut;

  \item the orientation of $\partial R$ and the base point $G(-1) = 0$, which together make the real
    segment contribute exactly zero;

  \item the running accumulation of $G$ across arcs -- $G$ is a single function on the whole
    boundary, not four independent ones.
\end{itemize}
A cheap self-test is available at no extra cost: after the full loop, $G$ must return to zero, and
the real part of $\oint f \overline{G} \, \dd Z$ must vanish.
Both hold to $7 \times 10^{-15}$ in our implementation and both catch orientation and branch errors
immediately.

For the four-tachyon contact term this replaces a two-dimensional quadrature by six boundary walks
per momentum point, through \eqref{jb:eq:pq-mandelstam}.
The gain is not only speed: it also removes the $n^{-3}$ convergence of the surface integral, which
is the accuracy bottleneck of the direct method at its usual settings.

What this does \emph{not} give is a closed form.
For generic $(p, q)$ we found none by the boundary route, and \cref{jb:sec:obstruction} says precisely
what stands in the way; on the two lines $p = 0$ and $q = 0$ a closed form in ${}_2F_1$ and Appell
$F_1$ is given with all its coefficients in \cref{jb:sec:obstruction:closure} and checked to
$10^{-15}$, and at the point $p = q = -1$ one in ${}_2F_1$ and ${}_3F_2$.
If a closed form is wanted for the general case, the two openings we would pursue are the series of
\cref{jb:sec:confocal}, whose only hard piece is the triangle adjacent to the unit-circle arc, and the
contour deformation of \cref{jb:sec:obstruction:lead}.




\subsection{The contact term depends on the momenta only through \texorpdfstring{$(s, t, u)$}{(s, t, u)}}
\label{sec:sl2c:stu}

\status{Conditional(PH29, PH30) -- Verified}\record{044-stu-reduction-contact-term (R48)}
The off-shell four-tachyon contact term of the declared section of the vertex family depends on the external momenta only through the Mandelstam invariants: the individual virtualities $k_a^2$ cancel pointwise in the modulus $Z$ and on each tile, using only momentum conservation and no mass-shell condition.
This is what makes the level-zero pair kernel of \cref{ch:momentum-space} a function of two invariants at fixed total momentum.
The derivation is exact; it was checked as a symbolic identity and numerically against an implementation of the local coordinates of the family.

\paragraph{Setup and inputs.}
All four momenta are incoming, $\sum_a k_a = 0$, and $s_{ab} := (k_a + k_b)^2$, so that
\begin{equation}
  k_a \cdot k_b
    = \frac{1}{2} \, \big( s_{ab} - k_a^2 - k_b^2 \big).
\end{equation}
The punctures sit at $(z_1, z_2, z_3, z_4) = (0, 1, \infty, Z)$.
The puncture at infinity is regularized by dropping every factor that carries the index $3$ as a puncture position, both in the mapping radii and in the Koba--Nielsen factor; the two removals must be the same prescription.
The tachyon vertex operator has conformal weight $(\frac{\alpha' k^2}{4} - 1, \frac{\alpha' k^2}{4} - 1)$, so that the local-frame dressing of puncture $a$ is $\rho_a^{\alpha' k_a^2 / 2 - 2}$, with $\rho_a$ the mapping radius.
The two inputs from the construction of the family~\cite{Erbin:2026:SL2CVerticesClosed} are the following.
\begin{enumerate}
  \item \emph{Closed form of the mapping radii.} On the declared section,
    \begin{equation}
      \label{eq:sl2c:radii}
      \rho_a
        = c(Z) \, P_a,
      \qquad
      c(Z)
        = \frac{1}{\beta \, (\abs{Z} \abs{1 - Z})^{1/2}},
      \qquad
      P_a
        := \sideset{}{'}\prod_{b \ne a} \abs{z_{ab}},
    \end{equation}
    where the prime removes the factors with the index $3$, and $c$ is the same at the four punctures.
  \item \emph{Assembly from twelve tiles.} With $G = S_3 \times \{ 1, \mathrm{conj} \}$, of order $12$,
    \begin{equation}
      \label{eq:sl2c:tiles}
      \{ T^4 \}_{\rm contact}
        = C \sum_{g \in G} \int_R \bigg[ \prod_a \rho_a(Z)^{\alpha' k_{\sigma_g^{-1}(a)}^2 / 2 - 2} \bigg] \, \mathrm{KN}(\sigma_g^{-1} k; Z) \, \dd^2 Z,
      \qquad
      \mathrm{KN}(k; Z)
        = \abs{Z}^{\alpha' k_1 \cdot k_4} \abs{1 - Z}^{\alpha' k_2 \cdot k_4},
    \end{equation}
    where $\sigma_g$ is the permutation of the momentum labels attached to the tile $g$ and $C$ is the overall constant of the vertex.
\end{enumerate}
The rest of the derivation uses nothing else.

\paragraph{Pointwise cancellation.}
With the punctures above, the regularized products are
\begin{equation}
  P_1
    = \abs{Z},
  \qquad
  P_2
    = \abs{1 - Z},
  \qquad
  P_3
    = 1,
  \qquad
  P_4
    = \abs{Z} \abs{1 - Z}.
\end{equation}
The momentum-dependent part of the dressing is then
\begin{equation}
  D(k; Z)
    := \prod_a \rho_a^{\alpha' k_a^2 / 2}
    = c(Z)^{\alpha' \sum_a k_a^2 / 2} \, \abs{Z}^{\alpha' (k_1^2 + k_4^2) / 2} \, \abs{1 - Z}^{\alpha' (k_2^2 + k_4^2) / 2},
\end{equation}
while the Koba--Nielsen factor reads, with the identity for $k_a \cdot k_b$,
\begin{equation}
  \mathrm{KN}(k; Z)
    = \abs{Z}^{\alpha' (s_{14} - k_1^2 - k_4^2) / 2} \, \abs{1 - Z}^{\alpha' (s_{24} - k_2^2 - k_4^2) / 2}.
\end{equation}
In the product the virtualities cancel exponent by exponent, and identically in $Z$,
\begin{equation}
  \label{eq:sl2c:cancel}
  D(k; Z) \, \mathrm{KN}(k; Z)
    = c(Z)^{\alpha' \sum_a k_a^2 / 2} \, \abs{Z}^{\alpha' s_{14} / 2} \, \abs{1 - Z}^{\alpha' s_{24} / 2}.
\end{equation}
The mechanism is structural: each $P_a$ is built from the same factors $\abs{z_{ab}}$ that appear in the Koba--Nielsen product, and $k_a \cdot k_b$ contains $- \frac{1}{2} (k_a^2 + k_b^2)$, so the dressing carries each $k_a^2$ with the opposite power.
The cancellation leaves a single power of $c$ only because $c$ does not depend on the puncture.
The remaining dependence on the virtualities is through their sum, which momentum conservation fixes:
\begin{equation}
  \label{eq:sl2c:sum}
  s + t + u
    = (k_1 + k_2)^2 + (k_1 + k_3)^2 + (k_1 + k_4)^2
    = 3 k_1^2 + k_2^2 + k_3^2 + k_4^2 + 2 k_1 \cdot (k_2 + k_3 + k_4)
    = \sum_a k_a^2,
\end{equation}
since $k_2 + k_3 + k_4 = - k_1$.
This is the off-shell form of the familiar on-shell identity; on the four-tachyon shell it gives $\alpha' (s + t + u) = 16$.
The momentum-independent part of the dressing is constant on the region:
\begin{equation}
  \label{eq:sl2c:product}
  \prod_a \rho_a
    = c(Z)^4 \prod_a P_a
    = \beta^{-4} \, (\abs{Z} \abs{1 - Z})^{-2} \, \abs{Z}^2 \abs{1 - Z}^2
    = \beta^{-4},
\end{equation}
so that $\prod_a \rho_a^{-2} = \beta^8$.
On the tile that carries the ordered pair $(x, y)$ of invariants in place of $(s_{14}, s_{24})$, the full integrand is therefore
\begin{equation}
  \label{eq:sl2c:tile-integrand}
  \mathcal I_{(x, y)}(s, t, u; Z)
    = \beta^8 \, c(Z)^{\alpha' (s + t + u) / 2} \, \abs{Z}^{\alpha' x / 2} \abs{1 - Z}^{\alpha' y / 2}.
\end{equation}

\paragraph{Tile sum and crossing.}
A tile does not move the integration region; it permutes the momentum labels.
Writing $(a, b, c, d) = (\sigma^{-1}(1), \sigma^{-1}(2), \sigma^{-1}(3), \sigma^{-1}(4))$, the computation above with $k \to \sigma^{-1} k$ gives~\eqref{eq:sl2c:tile-integrand} with $x = s_{ad}$ and $y = s_{bd}$.
Since $a, b, c$ are the three labels other than $d$, the three invariants $s_{ad}, s_{bd}, s_{cd}$ are $\{ s, t, u \}$ in some order: $(x, y)$ is an ordered pair of distinct invariants, and $z := s_{cd}$ is the third.
The permutations act on the cross-ratio with kernel the Klein group $V = \{ e, (12)(34), (13)(24), (14)(23) \}$, which fixes $s_{12} = s_{34}$, $s_{13} = s_{24}$ and $s_{14} = s_{23}$ each, and $S_4 / V \simeq S_3$ acts faithfully on $\{ s, t, u \}$.
The $S_3$ factor of $G$ is a set of coset representatives of $V$, so as $g$ runs over it, $(x, y)$ runs over the six ordered pairs once each.
The conjugation $Z \to \bar Z$ leaves~\eqref{eq:sl2c:tile-integrand} invariant and the momenta untouched, and contributes a factor $2$.
Inserting $c(Z)^{\alpha' (s + t + u) / 2} = \beta^{- \alpha' (s + t + u) / 2} (\abs{Z} \abs{1 - Z})^{- \alpha' (x + y + z) / 4}$, the exponents of $\abs{Z}$ and $\abs{1 - Z}$ become
\begin{equation}
  \frac{\alpha' x}{2} - \frac{\alpha' (x + y + z)}{4}
    = \frac{\alpha'}{4} (x - y - z),
  \qquad
  \frac{\alpha' y}{2} - \frac{\alpha' (x + y + z)}{4}
    = \frac{\alpha'}{4} (y - x - z),
\end{equation}
and the contact term is
\begin{equation}
  \label{eq:sl2c:contact}
  \{ T^4 \}_{\rm contact}(s, t, u)
    = 2 C \, \beta^{8 - \alpha' (s + t + u) / 2} \sum_{(x, y, z) \in S_3 \cdot (s, t, u)} J\Big( \frac{\alpha'}{4} (x - y - z), \frac{\alpha'}{4} (y - x - z) \Big),
\end{equation}
which is~\eqref{eq:mom:F} with $C$ restored, and the exponents~\eqref{jb:eq:pq-mandelstam} of the note.
The sum over the six orderings of a fixed function of $(x, y)$ is totally symmetric in $(s, t, u)$: given that the assembly~\eqref{eq:sl2c:tiles} sums over a transversal of $V$, total crossing symmetry is derived, not assumed separately.
The region is compact and bounded away from $Z = 0$ and $Z = 1$ for $\beta > \beta_{04}$, so each $J$ is finite for all complex exponents and the contact term is an entire function of the momenta.

\paragraph{Limits.}
At equal virtualities $k_a^2 = k^2$ the dressing factorizes before any use of the cancellation, $\prod_a \rho_a^{\alpha' k^2 / 2 - 2} = \beta^{8 - 2 \alpha' k^2}$ by~\eqref{eq:sl2c:product}, and~\eqref{eq:sl2c:contact} agrees since $s + t + u = 4 k^2$ on that line; the residual exponents match the equal-mass Koba--Nielsen exponents term by term, for example $\frac{\alpha'}{4} (s_{14} - s_{24} - s_{12}) = \frac{\alpha'}{2} s_{14} - \alpha' k^2 = \alpha' k_1 \cdot k_4$.
The equal-mass line cannot distinguish a function of $(s, t, u)$ from a function of the four virtualities, because on it the invariants fix $k^2$.
At $k = 0$ the contact term reduces to $12 \, C \, \beta^8 \Vol(R)$.

\paragraph{What the reduction uses.}
Exactly one property of the section is load-bearing: that $c$ in~\eqref{eq:sl2c:radii} is the same at every puncture.
With $\rho_a = c_a(Z) P_a$ and $c_a = c \, \e^{\gamma_a}$,~\eqref{eq:sl2c:cancel} acquires the factor
\begin{equation}
  \prod_a c_a(Z)^{\alpha' k_a^2 / 2}
    = c(Z)^{\alpha' (s + t + u) / 2} \exp\Big[ \frac{\alpha'}{2} \sum_a \gamma_a(Z) \, k_a^2 \Big],
\end{equation}
which is a function of $(s, t, u)$ alone if and only if the $\gamma_a$ coincide.
The reduction is thus an \emph{if and only if} statement about the section.
It survives any change of the second-point map $S$ of the section, which enters the radii only through the common factor $c = \abs{S}$; the constant $\beta^8$ of~\eqref{eq:sl2c:product}, by contrast, is special to the declared map $S_0$.
The weight of the interior point map of the section does not enter the radii, so the reduction does not depend on it.
A different consistent regularization of the puncture at infinity, for example a common finite cutoff, changes $c$ by a factor common to all punctures and leaves the reduction intact.

\paragraph{A residual weight on $\rho_3$.}
Two deformations of the radii by a momentum-independent constant $\mu > 0$ behave differently.
If only the radius at infinity is rescaled, $\rho_3 \to \mu \rho_3$, the tile $g$ acquires $\mu^{\alpha' k_{c_g}^2 / 2 - 2}$ with $c_g = \sigma_g^{-1}(3)$, and the deformed sum
\begin{equation}
  \{ T^4 \}^{\mu}_{\rm contact}
    = \mu^{-2} \, 2 C \sum_{(x, y)} \mu^{\alpha' k_{c(x, y)}^2 / 2} \int_R \mathcal I_{(x, y)}(s, t, u; Z) \, \dd^2 Z
\end{equation}
depends on the individual virtualities: the reduction to $(s, t, u)$ detects such a weight.
If all four radii are rescaled, $c \to \mu c$, the whole contact term is multiplied by
\begin{equation}
  \label{eq:sl2c:common-weight}
  \mu^{\alpha' (s + t + u) / 2 - 8},
\end{equation}
by~\eqref{eq:sl2c:sum}: the result is still a crossing-symmetric function of $(s, t, u)$, and no test of the momentum structure can see the weight.
On the four-tachyon shell, $\alpha' (s + t + u) = 16$, the exponent of~\eqref{eq:sl2c:common-weight} vanishes and the weight drops out for every $\mu$, so a test of this deformation must be run off that plane; the numerical check uses $\alpha' (s + t + u) = 6$.
Whether the first deformation shows up as a violation of crossing depends on which momentum each tile places at infinity, which the assembly~\eqref{eq:sl2c:tiles} does not fix (the label $\sigma^{-1}(3)$ takes all four values across a coset of $V$); what does not depend on that choice is that the weight destroys the reduction.

The implementation of the local coordinates against which the reduction was checked defines the map at infinity through the Klein relation $F_3 = Z / F_1$, so that its two routes to $\rho_3$ share that definition, and the check cannot by itself exclude a weight folded into it.
The next section closes this point.

\paragraph{Checks.}
The cancellation~\eqref{eq:sl2c:cancel} was verified as an exact symbolic identity,~\eqref{eq:sl2c:sum} as a polynomial identity in the four momenta, and~\eqref{eq:sl2c:product} numerically against radii obtained by differentiating the local coordinate maps, at both $\beta = \frac{3}{2}$ and $\beta = 2$.
The decisive test compares the integrand at two momentum configurations with the same invariants and different virtualities, in $D = 26$ with $\alpha' = 1$:
\begin{center}
  \small
  \begin{tabular}{lccccccc}
    \toprule
    & $\alpha' k_1^2$ & $\alpha' k_2^2$ & $\alpha' k_3^2$ & $\alpha' k_4^2$ & $\alpha' s_{12}$ & $\alpha' s_{13}$ & $\alpha' s_{14}$ \\
    \midrule
    X & $4$ & $4$ & $4$ & $4$ & $2$ & $5$ & $9$ \\
    Y & $-1$ & $7$ & $3$ & $7$ & $2$ & $5$ & $9$ \\
    X$'$ & $1$ & $1$ & $2$ & $2$ & $-4$ & $5$ & $5$ \\
    Y$'$ & $-1$ & $3$ & $0$ & $4$ & $-4$ & $5$ & $5$ \\
    \bottomrule
  \end{tabular}
\end{center}
% NOTE: the derivation names the invariants t = s_13, u = s_14, the opposite of the note's convention; the columns are labelled by s_ab to avoid the clash.
The configuration X is the equal-mass on-shell point, Y shares its invariants with every leg off shell, and the second pair lies off the plane $\alpha' (s + t + u) = 16$; all four were realized as explicit vectors in $\R^{1,25}$.
The dressing and the Koba--Nielsen factor differ between X and Y pointwise in $Z$, only their product agrees, which is why the test is not vacuous.
The integrands agree pointwise and after quadrature; deformations that break the cancellation, momentum conservation, the identity~\eqref{eq:sl2c:product}, the assignment of $t$ and $u$, and the radius at infinity were each confirmed to make the corresponding check fail.

\paragraph{Scope.}
The statement concerns four tachyons and the contact term only.
For an excited external state the dressing exponent shifts with the level and the phase of the local frame enters, which the modulus $\rho_a$ does not see; the exchange contributions carry radii specific to the Feynman regions and are outside the assembly~\eqref{eq:sl2c:tiles}.
The statement is about the integrand, not about the numerical tractability of the resulting function, and it does not establish the two inputs, which it inherits.

\paragraph{Consequence for the pair kernel.}
The level-zero pair kernel has external momenta $k_1 = k$, $k_2 = P - k$, $k_3 = - k'$, $k_4 = - (P - k')$, with virtualities that differ in general, which suggests a function of six parameters.
By the reduction it depends on $(k, k')$ only through $(s, t, u)$, and $s = P^2$ is fixed: at fixed total momentum the kernel is a function of the two invariants $t$ and $u$, with $t + u = \sum_a k_a^2 - P^2$.
Total crossing symmetry includes $t \leftrightarrow u$, the exchange of the two strands of one pair, which gives an exact check on any discretization.
This is the form~\eqref{eq:mom:KP} used in \cref{sec:mom:level-0}.

\subsection{No residual factor on the mapping radius at infinity}
\label{sec:sl2c:rho3}

\record{052-rho3-residual-factor (R54)}
The radius at infinity carries no residual factor on the declared section, so the reduction of the previous section holds with no weight that depends on the tile.
Two statements, with different statuses, establish this.
\begin{itemize}
  \item For any assignment of local coordinates symmetric under $S_4$, $\rho_3 = \rho_1 / \abs{Z}$, and the plumbing of two cubic vertices gives $\tilde F_3 = Z / \tilde F_1$: the Klein relation is the $(13)(24)$ symmetry of the vertex, not an additional input. \status{Proven}
  \item $\abs{(1 / \tilde F_3)'(0)} = c(Z)$, the residual factor is $f \equiv 1$ on the declared section, and the reduction carries no tile-dependent weight on the twelve tiles. \status{Conditional(PH29, PH30) -- Proven}
    The identification $\rho_1 = \abs{S_0} = c \abs{Z}$ uses the closed form~\eqref{eq:sl2c:radii}, and the tiles of the $t$- and $u$-channels are reached through the relabelling of the assembly~\eqref{eq:sl2c:tiles}, not computed separately.
\end{itemize}

\paragraph{The question.}
Define on the region
\begin{equation}
  f(Z)
    := \frac{\abs{(1 / F_3)'(0)}}{c(Z)},
\end{equation}
the ratio of the radius at infinity, read in the chart $t = 1 / z$, to its closed form.
A nontrivial $f$ would multiply the integrand of the tile $\sigma$ by $f(Z)^{\alpha' k_{\sigma^{-1}(3)}^2 / 2 - 2}$.
Each of $k_1, k_2, k_3$ sits at infinity on two tiles, so the weight is the same on all tiles only when $k_1^2 = k_2^2 = k_3^2$; otherwise the virtualities no longer cancel and the reduction fails.
The maps $F_a$ carry no momentum, so $f$ depends only on $(Z, \beta)$ and one scan over the region decides the question for every assignment of masses.

\paragraph{First route: symmetry of the vertex.}
The element $(13)(24)$ of $V$ is realized by the Möbius map $M_{13}(z) = Z / z$, which exchanges $0 \leftrightarrow \infty$ and $1 \leftrightarrow Z$ and fixes the cross-ratio.
Invariance of the vertex under this relabelling relates the local coordinates at the punctures $1$ and $3$ up to a rotation of the local coordinate, $M_{13} \circ F_1(W, Z) = F_3(\e^{\I \theta} W, Z)$, so that
\begin{equation}
  \frac{1}{F_3(W')}
    = \frac{F_1(\e^{- \I \theta} W')}{Z},
  \qquad
  \abs{\Big( \frac{1}{F_3} \Big)'(0)}
    = \frac{\abs{F_1'(0)}}{\abs{Z}}
    = \frac{\abs{S}}{\abs{Z}}.
\end{equation}
The rotation $W \to \e^{\I \theta} W$, which a closed-string vertex allows at each puncture because of level matching, drops out of the modulus of the derivative.
Hence $\rho_3 = \rho_1 / \abs{Z}$ for every $S_4$-symmetric choice of local coordinates, whatever the maps $S$ and $R$ of the section.
With the declared map, $\abs{S_0}^2 = \abs{Z} / (\beta^2 \abs{1 - Z})$, and
\begin{equation}
  \rho_3
    = \frac{1}{\beta \, \abs{Z}^{1/2} \abs{1 - Z}^{1/2}}
    = c(Z).
\end{equation}
A residual factor on $\rho_3$ would therefore not be a property of a particular section: it would be a violation of the $(13)(24)$ symmetry of the vertex.

\paragraph{Second route: plumbing of two cubic vertices.}
The same relation follows from the plumbing construction of the $s$-channel maps, without using the Klein relation.
In the cubic vertex of the family the local coordinate at the puncture at infinity is $w_\infty = 2 \beta / (2 z - 1)$, so that the plumbing $w^{(1)}_\infty w^{(2)}_\infty = q$ reads $(2 z^{(1)} - 1)(2 z^{(2)} - 1) = 4 \beta^2 / q$.
With the assignment $W_1 = w^{(1)}_0$, $W_2 = w^{(1)}_1$, $W_3 = w^{(2)}_0$, $W_4 = w^{(2)}_1$, and the Möbius map $g_s(z) = (p_3 - 1) z / (p_3 - z)$ that sends $p_3 = F_3^{\rm pl}(0)$ to infinity and fixes $0$ and $1$,
\begin{equation}
  F_3^{\rm pl}(W)
    = \frac{1}{2} \Big( 1 + \frac{4 \beta^2}{q} \, \frac{1}{2 f_0(W) - 1} \Big),
  \qquad
  \tilde F_a
    = g_s \circ F_a^{\rm pl},
  \qquad
  Z
    = \tilde F_4(0)
    = \frac{(4 \beta^2 + q)^2}{16 \beta^2 q},
\end{equation}
where $f_0$ is the local coordinate map at $0$ of the cubic vertex.
The relations
\begin{equation}
  \tilde F_3 - \frac{Z}{\tilde F_1}
    = 0,
  \qquad
  \tilde F_4 - \frac{Z}{\tilde F_2}
    = 0
\end{equation}
then hold as rational identities in $(W, q, \beta)$, and
\begin{equation}
  \Big( \frac{1}{\tilde F_3} \Big)'(0)
    = \frac{16 \beta q}{(4 \beta^2 - q)(4 \beta^2 + q)}.
\end{equation}
Since $1 - Z = - (4 \beta^2 - q)^2 / (16 \beta^2 q)$,
\begin{equation}
  \abs{Z} \abs{1 - Z}
    = \frac{\abs{4 \beta^2 + q}^2 \abs{4 \beta^2 - q}^2}{256 \beta^4 \abs{q}^2},
  \qquad
  c(Z)
    = \frac{16 \beta \abs{q}}{\abs{4 \beta^2 + q} \abs{4 \beta^2 - q}}
    = \abs{\Big( \frac{1}{\tilde F_3} \Big)'(0)},
\end{equation}
for every $q$, in particular for the root $q(Z)$ on the region.
The Klein relation used by the implementation is therefore exactly what the plumbing gives.

\paragraph{Checks.}
The plumbing identities were verified symbolically, and the derivative against $c(Z)$ to $1.2 \times 10^{-14}$ at $80$ points of the region.
A scan of the implemented maps against the closed form, on $250$ points of the region and $16$ boundary points for each $\beta$, gives $\max \abs{f - 1} = 6.4 \times 10^{-14}$ at $\beta = \frac{3}{2}$ and $1.7 \times 10^{-14}$ at $\beta = 2$; the residual decreases as $h^4$ with the step of the fourth-order difference stencil, down to a plateau at $1.2 \times 10^{-15}$.
A test of the way $\rho_3$ enters the integrand was also run: moving the four maps by a random Möbius transformation, so that all punctures are finite and no regularization is needed, reproduces the regularized tile integrand at unequal masses to $1.1 \times 10^{-11}$, and fails by more than $7.7 \times 10^{-2}$ when $\rho_3$ is multiplied by $\abs{Z}$.
This last test checks the chart and the regularization of the puncture at infinity, not the Klein relation.
Deformations of the map at infinity, of the closed form, of the pairing of punctures and of $f$ were each confirmed to make the corresponding check fail; a deformation $F_3 (1 + \epsilon w)$ cannot be seen, since it is proportional to $1 / F_3(0) = 0$.

\paragraph{Scope.}
The first route uses only the $(13)(24)$ element of $V$, for any section symmetric under it; the declared section is one by its construction.
The statement concerns the mapping radii only: the dressing of excited states, which involves the interior point map and the higher Taylor coefficients of the local coordinates, is not covered.

\subsection{Corner asymptotics of the tile integral}
\label{sec:sl2c:corner}

\record{049-kp-pair-kernel (R52)}
For $q \to + \infty$ at fixed $\delta \ge 0$,
\begin{equation}
  \label{eq:sl2c:corner}
  J(- q - \delta, q)
    = 2 \sqrt{2 \pi} \; 2^q \, q^{-3/2} \, \big[ 1 + O(1 / q) \big].
\end{equation}
\status{Proven} at leading order; the coefficient of the $1 / q$ correction is observed numerically, not derived.
Its consequence for the decay of the diagonal of the pair kernel is \status{Conditional(PH25, PH29, PH30) -- Numerical}: it holds under the Euclidean prescription of \cref{sec:mom:level-0} and the two inputs of \cref{sec:sl2c:stu}.
In this section $p$ and $q$ are the exponents of $J$.

\paragraph{Where the integral concentrates.}
For $p + q \le 0$ and $q$ large, $\abs{Z}^p \abs{1 - Z}^q$ is largest where $\abs{1 - Z}$ is large and $\abs{Z}$ small, which on $R$ is the corner $Z = - 1$, where the unit circle and the real segment of $\partial R$ meet; the region lies outside the unit circle there.
Write $Z = - 1 + \epsilon$ with $\epsilon = \epsilon_1 + \I \epsilon_2$ and $\epsilon_2 \ge 0$.
The constraint $\abs{Z} \ge 1$ reads $\epsilon_1 \le \frac{1}{2} (\epsilon_1^2 + \epsilon_2^2)$, that is $\epsilon_1 \le \frac{1}{2} \epsilon_2^2 + O(\epsilon_2^4)$.
Expanding the logarithms,
\begin{equation}
  \ln \abs{Z}
    = - \epsilon_1 - \frac{1}{2} (\epsilon_1^2 - \epsilon_2^2) + O(\abs{\epsilon}^3),
  \qquad
  \ln \abs{1 - Z}
    = \ln 2 - \frac{\epsilon_1}{2} - \frac{1}{8} (\epsilon_1^2 - \epsilon_2^2) + O(\abs{\epsilon}^3),
\end{equation}
the exponent with $p = - q - \delta$ is
\begin{equation}
  \Phi
    := p \ln \abs{Z} + q \ln \abs{1 - Z}
    = q \ln 2
      + \epsilon_1 \Big( \frac{q}{2} + \delta \Big)
      + (\epsilon_1^2 - \epsilon_2^2) \Big( \frac{3 q}{8} + \frac{\delta}{2} \Big)
      + O(q \abs{\epsilon}^3).
\end{equation}
$\Phi$ increases with $\epsilon_1$, so near the corner its maximum over the region lies on the circle $\epsilon_1 = \frac{1}{2} \epsilon_2^2$, where
\begin{equation}
  \Phi
    = q \ln 2 - \frac{q}{8} \, \epsilon_2^2 + O(q \epsilon_2^4).
\end{equation}
The corner is a strict maximum, of width $\epsilon_2 \sim q^{-1/2}$ along the circle and $\xi := \epsilon_1 - \frac{1}{2} \epsilon_2^2 \sim - 1 / q$ across it.

\paragraph{Laplace's method.}
In the variables $(\xi, \epsilon_2)$, whose Jacobian is $1$, the terms in $\xi^2$ and $\xi \epsilon_2^2$ are of relative order $1 / q$ in this scaling, and
\begin{equation}
  J(- q - \delta, q)
    = 2^q \int_0^\infty \dd \epsilon_2 \, \e^{- q \epsilon_2^2 / 8} \int_{- \infty}^0 \dd \xi \, \e^{\xi (q / 2 + \delta)} \, \big[ 1 + O(1 / q) \big]
    = 2^q \cdot \frac{1}{q / 2 + \delta} \cdot \frac{1}{2} \sqrt{\frac{8 \pi}{q}} \, \big[ 1 + O(1 / q) \big],
\end{equation}
which is~\eqref{eq:sl2c:corner} at leading order: the exponent grows linearly across the unit circle, at rate $q / 2$, and quadratically along it, at rate $q / 8$, and the two widths give $q^{-1} \cdot q^{-1/2}$.
Writing the correction as $1 + c_1(\delta) / q$, the numerical values at $q = 200, 400, 800$ are $c_1(0) \approx 4.8$ and $c_1(8/3) \approx - 0.58$, consistent with $c_1(\delta) = c_1(0) - 2 \delta$, the $- 2 \delta$ being the expansion of $1 / (q / 2 + \delta)$.
The corner, and with it the leading asymptotics, does not depend on $\beta$: $J(-250, 250)$ agrees at $\beta = \frac{3}{2}$ and $\beta = 2$ to twelve digits, the rest of the region being exponentially suppressed.
The check compares $J \, q^{3/2} / (2 \sqrt{2 \pi} \, 2^q) - 1$ with the bound $8 / q$ at the three values of $q$ and $\delta \in \{ 0, \frac{8}{3} \}$; the same criterion fails when the power $q^{-3/2}$ is replaced by $q^{-1}$.

\paragraph{Consequence for the pair kernel.}
On the diagonal of the level-zero kernel, at relative momenta of equal modulus $\rho$ and large $\alpha' \rho^2$, the tiles that dominate have $q$ of order $\alpha' \rho^2$, and the corner power $q^{-3/2}$ steepens the decay.
The logarithmic slope of the symmetrized s-wave diagonal in $\alpha' \rho^2$ lies below the edge rate $- (2 \ln \beta - \ln 2)$ by $\frac{3}{2} / (\alpha' \rho^2)$, partly compensated by the angular average, which pulls the slope towards the edge rate.
For one of the two dominant tiles the exponent $\delta$ is itself of order $q$ on the relevant angular window, and its corner factor is $2 \sqrt{2 \pi} \, 2^q q^{-1/2} (q / 2 + \delta)^{-1}$ rather than~\eqref{eq:sl2c:corner}; both corner factors are of order $\rho^{-3}$.
At $\beta = \frac{3}{2}$ the measured slope is $- 0.123 \pm 0.001$ on $\alpha' \rho^2 \in [60, 250]$, against the edge rate $- 0.1178$; the model with the two corner factors and the exact angular average reproduces the five measured slopes to $1.05 \times 10^{-3}$, while a model with the angular average alone misses by $2.4 \times 10^{-2}$, on the wrong side.
The kernel itself and its spectral properties are treated in \cref{sec:mom:level-0}.

\subsection{Conjugation invariance of the contact region}
\label{sec:sl2c:conj}

\record{054-ph6-conjugation-invariance-n4 (R56)}
The reality of the quartic vertex, $\vertex{4}(A_1, \ldots, A_4)^* = \vertex{4}(A_1^\ddagger, \ldots, A_4^\ddagger)$, follows from the reality of the conformal field theory correlators~\cite{Sonoda:1990:CovariantClosedString} together with the invariance of the vertex region and of its measure under complex conjugation of the moduli (\cref{ch:conventions}).
This section checks the geometric inputs for the declared section of the family at four points, for $\beta \ge \beta_{04}$.
Its statements have different statuses.
\begin{itemize}
  \item The conjugation invariance of the region, the consistency of its tiling by $G = S_3 \times \{ 1, \mathrm{conj} \}$, and the conjugation equivariance of the section data to first order: \status{Proven}.
  \item The reality of the level-zero four-tachyon contact term at real off-shell momenta: \status{Conditional(PH29, PH30) -- Proven}.
  \item The reality of each pair of conjugate tiles for primaries of integer spin: \status{Proven}, given the reality of plane correlators with real exponents taken from~\cite{Sonoda:1990:CovariantClosedString}.
  \item That the measure $2$-form with its $b \bar b$ insertion is real and invariant under conjugation up to orientation beyond level zero: \status{Argued}.
\end{itemize}
That the region defined below is the contact region of the family is taken from its construction~\cite{Erbin:2023:SL2CQuarticVertex, Erbin:2026:SL2CVerticesClosed}.

\paragraph{The region.}
The fundamental domain is $\mathrm{FD} = \{ \Im Z \ge 0, \ \abs{Z} \ge 1, \ \Re Z \le \frac{1}{2} \}$, with walls $C_1$ (the real half-line $Z \le - 1$), $C_2$ ($\Re Z = \frac{1}{2}$) and $C_3$ ($\abs{Z} = 1$).
$E_\beta$ is the closed ellipse centred at $\frac{1}{2}$ with semi-axes $a$ and $b$ of~\eqref{jb:eq:axes} along the real and imaginary axes; its boundary $C_4$ is the image of $\abs{q} = 1$ under $Z^{(s)}(q) = \frac{1}{2} + \beta^2 / q + q / 16 \beta^2$.
The tile is $T := \mathrm{FD} \cap E_\beta$, and $S_3$ is the anharmonic group $\{ Z, 1 - Z, 1 / Z, Z / (Z - 1), 1 / (1 - Z), (Z - 1) / Z \}$.
The $s$-channel Feynman region is the exterior of $E_\beta$, and the $t$- and $u$-channel regions are its images under the maps that send infinity to $0$ and to $1$, so the contact region has the definition, independent of the tiling,
\begin{equation}
  \label{eq:sl2c:region-ind}
  \mc V_{0,4}
    = \big\{ Z : \ Z \in E_\beta, \ 1 / Z \in E_\beta, \ 1 / (1 - Z) \in E_\beta \big\}.
\end{equation}

\paragraph{Invariance and tiling.}
Every element of $S_3$ is a Möbius map with integer coefficients, so $\overline{h(Z)} = h(\bar Z)$ and
\begin{equation}
  \mathrm{conj} \circ h
    = h \circ \mathrm{conj},
  \qquad
  h \in S_3:
\end{equation}
$G$ is a direct product of order $12$, and conjugation permutes the tiles, $\mathrm{conj}(h(T)) = (h \cdot \mathrm{conj})(T)$.
The ellipse is symmetric under $Z \to \bar Z$, since its centre is real and its axes are along the real and imaginary directions, so each condition of~\eqref{eq:sl2c:region-ind} is invariant and $\overline{\mc V_{0,4}} = \mc V_{0,4}$.
The domain $\mathrm{FD}$ is one of the twelve triangles cut out of the sphere by the real axis, the circles $\abs{Z} = 1$ and $\abs{Z - 1} = 1$ and the line $\Re Z = \frac{1}{2}$, with vertices $-1$, $\e^{\I \pi / 3}$ and $\infty$; these triangles are the orbit of $\mathrm{FD}$ under $G$, which acts simply transitively on them.
Hence $\mathrm{conj}(T)$ is the tile of $g = \mathrm{conj}$, and $T \cap \mathrm{conj}(T) = C_1 \cap E_\beta$ has measure zero.

The one input specific to the family is that the twelve tiles reproduce~\eqref{eq:sl2c:region-ind}.
The conditions of~\eqref{eq:sl2c:region-ind} are permuted by $S_3$ and invariant under conjugation, so the region is invariant under $G$, and the statement reduces to $\mathrm{FD} \cap \mc V_{0,4} = \mathrm{FD} \cap E_\beta$.
On $\mathrm{FD}$, $\abs{Z} \ge 1$ and $\abs{1 - Z} \ge 1$, so $1 / Z$ and $1 / (1 - Z)$ lie in the closed unit disc, and the reduction holds if and only if the unit disc lies in $E_\beta$.
The squared distance from $0$ to the point $(\frac{1}{2} + a \cos \theta, b \sin \theta)$ of the ellipse is, with $u = \cos \theta$ and $c = \frac{1}{2}$,
\begin{equation}
  d^2(\theta)
    = (a^2 - b^2) \, u^2 + 2 a c \, u + c^2 + b^2,
\end{equation}
a convex quadratic in $u$, since $a^2 - b^2 = \frac{1}{4}$, with vertex at $u^* = - a c / (a^2 - b^2) = - 2 a \le - 1$, since $a \ge \frac{1}{2}$.
Its minimum on $[-1, 1]$ is at $u = -1$, the left vertex, where $d^2 = (a - c)^2$, so that
\begin{equation}
  \operatorname{dist}(0, \partial E_\beta)
    = a - \frac{1}{2}
\end{equation}
exactly, for every $\beta > \frac{1}{2}$.
The unit disc lies in $E_\beta$ if and only if $a \ge \frac{3}{2}$, that is $16 \beta^4 - 24 \beta^2 + 1 \ge 0$, that is $\beta \ge \beta_{04}$; for $\beta = \frac{3}{2}$ and $2$ the distance is $1.78$ and $3.52$.
On the boundary, conjugation fixes $C_1$ pointwise and maps $C_2$, $C_3$, $C_4$ onto the corresponding arcs of $\partial(\mathrm{conj} \, T)$; with the Schwarz functions of the arcs (\cref{jb:sec:obstruction:schwarz}), each interior wall is the fixed set of an orientation-reversing element $\mathrm{conj} \circ h$.
The group $G$ is generated by the reflections in the three walls of $\mathrm{FD}$, a spherical triangle group with angles $\frac{\pi}{2}$, $\frac{\pi}{3}$, $\frac{\pi}{2}$ at $-1$, $\e^{\I \pi / 3}$, $\infty$, and conjugation is the reflection in $C_1$.
The region is therefore conjugation invariant, and conjugation permutes the twelve tiles in six pairs $\{ h(T), h(\bar T) \}$.

\paragraph{The section data.}
The data of the section are functions of $Z$ with real coefficients, or principal branches whose cuts lie on the real axis, and for such functions $f(\bar u) = \overline{f(u)}$ off the cut.
The plumbing parameter $q(Z)$ is the root with $\abs{q} \le 1$ of $q^2 + 8 \beta^2 (1 - 2 Z) q + 16 \beta^4 = 0$, a quadratic with real coefficients; the maps $F_a$ are rational in $(w, q, \beta)$ with real coefficients, with $F_{3,4} = Z / F_{1,2}$; and $S_0 = \beta^{-1} \sqrt{Z / (Z - 1)}$ has its cut at $Z \in [0, 1)$, outside $\mathrm{FD} \cup \overline{\mathrm{FD}}$.
Hence
\begin{equation}
  q(\bar Z)
    = \overline{q(Z)},
  \qquad
  F_a(\bar w; \bar Z)
    = \overline{F_a(w; Z)},
  \qquad
  \rho_a(\bar Z)
    = \rho_a(Z),
  \qquad
  S_0(\bar Z)
    = \overline{S_0(Z)},
\end{equation}
with $F_a'(0) = \kappa_a S_0$, $\kappa = (1, (1 - Z) / Z, 1 / Z, Z - 1)$.
The interior point map $R$ of the section, which is not holomorphic in $Z$, is equivariant in the form $R(\bar Z, Z) = \overline{R(Z, \bar Z)}$, because its non-holomorphic dependence enters only through a real weight that is invariant under conjugation.
The section on the tile $\bar T$ is the transport by conjugation of the section on $T$, and by these relations it is given by the same closed forms evaluated at $\bar Z$.
It is continuous across the mirror $C_1$ if and only if the data are real there, which they are: for $Z \le - 1$, $Z / (Z - 1) \in (\frac{1}{2}, 1)$ and $q$, $S_0$, $R$ and $F_a(w)$ at real $w$ are real.
On the two other walls the transport uses $\mathrm{conj} \circ (1 - Z)$ and $\mathrm{conj} \circ (1 / Z)$, and the corresponding conditions are the wall conditions of the section, which were re-checked.

\paragraph{Measure and level zero.}
With $\dd^2 Z$ a positive density and $Z \to \bar Z$ an isometry of Jacobian $- 1$,
\begin{equation}
  \int_{\bar T} f(Z) \, \dd^2 Z
    = \int_T f(\bar Z) \, \dd^2 Z
\end{equation}
as densities.
The reversal of orientation is absorbed by the $b \bar b$ insertion of the measure $2$-form, which at level zero is folded into the $- 2$ of the dressing exponent and is a positive density; beyond level zero this step is argued, not derived.
The tile integrands of the assembly~\eqref{eq:sl2c:tiles} depend on $Z$ through $\abs{Z}$, $\abs{1 - Z}$ and the radii, so
\begin{equation}
  \mathcal I_g(k; \bar Z)
    = \mathcal I_g(k; Z),
  \qquad
  \mathcal I_g(\bar k; Z)
    = \overline{\mathcal I_g(k; Z)},
  \qquad
  \mathcal I_g(k; Z)
    > 0
  \quad
  (k \text{ real}),
\end{equation}
and the six pairs of conjugate tiles add:
\begin{equation}
  \{ T(k_1) \cdots T(k_4) \}
    = C \sum_{h \in S_3} \Big[ \int_T \mathcal I_h + \int_{\bar T} \mathcal I_{h \, \mathrm{conj}} \Big]
    = 2 C \sum_{h \in S_3} \int_T \mathcal I_h(k; Z) \, \dd^2 Z.
\end{equation}
For real momenta every term is positive and $C$ is real, so the block is real.
The tachyon state $T(k) = c_1 \bar c_1 \ket{k}$ has $T(k)^\ddagger = \varepsilon \, T(-k)$ for real $k$, with a fixed sign $\varepsilon = \pm 1$, and the integrand is even in $k$, so
\begin{equation}
  \{ T(k_1) \cdots T(k_4) \}^*
    = \{ T(k_1) \cdots T(k_4) \}
    = \varepsilon^4 \{ T(- k_1) \cdots T(- k_4) \}
    = \{ T(k_1)^\ddagger \cdots T(k_4)^\ddagger \}.
\end{equation}

The level-zero block is real tile by tile: it would stay real with a tile removed, or with the map $S_0$ deformed by an imaginary term, since only $\abs{S}$ enters the radii.
It therefore does not test the geometric inputs, and every level-zero number of this document is blind to them.
To make them load-bearing, the check adds a probe sensitive to the phases: primaries of integer spin and real weights, inserted in the frames $F_a'(0) = \kappa_a S_0$, times a plane correlator of the compact-boson form with real exponents.
For such states the reality of the correlator together with the equivariance of the frames makes each pair of conjugate tiles real; this is a corollary of the equivariance of the section data, not an independent test of the invariance of the region.
The probe is real with the declared section and fails when the conjugate tile is dropped or when $S_0$ is deformed by $\I \eta \abs{Z} / \beta$, $\eta$ real.

\paragraph{Checks.}
The commutation with conjugation is exact, the group closes with twelve elements, the tiling reproduces~\eqref{eq:sl2c:region-ind} at $11\,980$ and $18\,035$ sample points of the region (out of $40\,000$ per $\beta$) with no mismatch and no double covering, and the data equivariance holds to rounding at $400$ random points.
The checks fail as they should for a Möbius map with non-real coefficients, for $\beta = 0.9 < \beta_{04}$, with the conjugate tile removed, with the deformed $S_0$, and with a wrong assignment of the mirrors.

\paragraph{Scope.}
The statement is at four points.
Excited states need the local coordinate maps of the section beyond first order, which differ from the plumbing maps through $R$; their equivariance is expected, since the construction from $(S_0, R)$ uses wall conditions with real coefficients and a real weight, but it was not checked.
The ghost measure beyond level zero is argued, and the reality of the conformal field theory correlators is cited.
For five or more points nothing is established here.
